\documentclass[11pt]{article}
% Lineage: v6-arg-core + spec-annex
% Last structural merge: 2026-02-11

% ------------------------------------------------------------
% Core packages (inlined from former class dependency)
% ------------------------------------------------------------
\usepackage[T1]{fontenc}
\usepackage{lmodern}
\usepackage{amsmath,amssymb,amsfonts,amsthm}
\usepackage{microtype}
\usepackage{geometry}
\usepackage{xcolor}
\usepackage{graphicx}
\usepackage{titlesec}
\usepackage{tocloft}
\usepackage{enumitem}
\usepackage{booktabs}
\usepackage{chngcntr}
\usepackage{hyperref}
\geometry{margin=1in}
\definecolor{AIONBlue}{RGB}{20, 60, 120}
\definecolor{AIONAccent}{RGB}{120, 20, 120}
\hypersetup{
    colorlinks=true,
    linkcolor=AIONBlue,
    citecolor=AIONBlue,
    urlcolor=AIONAccent,
    hypertexnames=false
}
\titleformat{\section}{\large\bfseries\color{AIONBlue!85!black}}{\thesection}{0.5em}{}
\titleformat{\subsection}{\normalsize\bfseries\color{AIONBlue!85!black}}{\thesubsection}{0.5em}{}
\counterwithin{figure}{section}
\renewcommand{\theHfigure}{\thesection.\arabic{figure}}

\theoremstyle{definition}
\newtheorem{definition}{Definition}[section]
\newtheorem{assumption}[definition]{Assumption}
\theoremstyle{plain}
\newtheorem{proposition}[definition]{Proposition}
\newtheorem{theorem}[definition]{Theorem}
\newtheorem{lemma}[definition]{Lemma}
\newtheorem{corollary}[definition]{Corollary}
\theoremstyle{remark}
\newtheorem{remark}[definition]{Remark}
\newtheorem{example}[definition]{Example}

% ------------------------------------------------------------
% Metadata commands + title page (inlined from former class dependency)
% ------------------------------------------------------------
\newcommand{\papertitle}{}
\newcommand{\papernumber}{}
\newcommand{\paperdate}{}
\newcommand{\paperversion}{}
\newcommand{\paperauthor}{}
\newcommand{\paperaffiliation}{}
\newcommand{\paperorcid}{}
\newcommand{\paperemail}{}
\newcommand{\paperdoi}{}

\newcommand{\AIONTitlePage}{
    \thispagestyle{empty}
    \begin{center}
        \vspace*{1.5cm}
        {\Huge\bfseries \papertitle \par}
        \vspace{14pt}
        {\normalsize\scshape\color{AIONBlue}
        AI$\Omega$N Foundations Series --- \papernumber \par}
        \vspace{16pt}
        {\large \paperauthor \par}
        \vspace{4pt}
        {\normalsize
        \paperaffiliation \\
        ORCID: \paperorcid \\
        Email: \href{mailto:\paperemail}{\paperemail} \par}
        \vspace{10pt}
        \if\relax\detokenize\expandafter{\paperdoi}\relax\else
          {\normalsize DOI: \href{https://doi.org/\paperdoi}{\paperdoi}\par}
        \fi
        {\normalsize Version: \paperversion \par}
        {\normalsize \paperdate \par}
    \end{center}
    \vspace{2cm}
}

% ------------------------------------------------------------
% Metadata for this paper
% ------------------------------------------------------------
\renewcommand{\papertitle}{WARP Graphs: Ethics of Deterministic Replay \& Provenance Sovereignty}
\renewcommand{\papernumber}{Paper VI}
\renewcommand{\paperversion}{v1.0.0}
\renewcommand{\paperdate}{March 2026}
\renewcommand{\paperdoi}{10.5281/zenodo.18863648}

\renewcommand{\paperauthor}{James Ross}
\renewcommand{\paperaffiliation}{Independent Researcher}
\renewcommand{\paperorcid}{0009-0006-0025-7801}
\renewcommand{\paperemail}{james@flyingrobots.dev}

% ------------------------------------------------------------
% Extra packages needed by this paper
% ------------------------------------------------------------
\usepackage{float}
\usepackage{xurl}
\usepackage{tikz}
\usepackage{tabularx}
\usepackage{colortbl}
\usetikzlibrary{arrows.meta,positioning,decorations.pathreplacing,calc,shapes.geometric,fit}

% Inlined support macros/diagrams (previously in macros.tex + diagrams.tex)
\newcommand{\AION}{\textrm{AI}\ensuremath{\Omega}\textrm{N}}
\newcommand{\AIONProjectURL}{\url{https://github.com/flyingrobots/aion}}
\DeclareRobustCommand{\WARP}{\ifmmode\mathit{WARP}\else{\upshape\scshape warp}\fi}
\DeclareMathOperator{\In}{In}
\DeclareMathOperator{\Out}{Out}
\newcommand{\Replay}{\mathsf{Replay}}
\newcommand{\skel}{\mathsf{skel}}

\newcommand{\AIONDiagramHolography}{
\begin{tikzpicture}[
    state/.style={circle, draw, fill=gray!20, minimum size=8mm, inner sep=0pt},
    boundary/.style={draw, rounded corners, thick, dashed, inner sep=6pt}
  ]
  \node[state] (u0) at (0,0) {$U_0$};
  \node[state] (u1) at (2,0) {$U_1$};
  \node[state] (u2) at (4,0) {$U_2$};
  \node[state] (un) at (6,0) {$U_n$};

  \draw[->, thick] (u0) -- (u1) node[midway, above] {\scriptsize $\mu_0$};
  \draw[->, thick] (u1) -- (u2) node[midway, above] {\scriptsize $\mu_1$};
  \draw[->, thick] (u2) -- (un) node[midway, above] {\scriptsize $\cdots$};

  \node[boundary, fit=(u0)] (b0) {};
  \node[below=2mm of b0] {\scriptsize $U_0$};

  \node[boundary, fit=(u0) (u1) (u2)] (bp) {};
  \node[below=2mm of bp] {\scriptsize payload $P$};

  \node[boundary, fit=(u0) (bp)] (b) {};
  \node[above=2mm of b] {\scriptsize boundary $(U_0,P)$};

  \draw[->, thick, dashed] (b.east) to[out=0,in=180] (un.west);
\end{tikzpicture}
}

\newcommand{\AIONDiagramWormhole}{
\begin{tikzpicture}[
    state/.style={circle, draw, fill=gray!20, minimum size=8mm, inner sep=0pt},
    wormhole/.style={->, thick, dashed},
    normal/.style={->, thick}
  ]
  \node[state] (u0) at (0,0) {$U_i$};
  \node[state] (u1) at (2,0) {$U_{i+1}$};
  \node[state] (u2) at (4,0) {$U_{i+2}$};
  \node[state] (uk) at (6,0) {$U_{i+k}$};

  \draw[normal] (u0) -- (u1) node[midway, above] {\scriptsize $\mu_i$};
  \draw[normal] (u1) -- (u2) node[midway, above] {\scriptsize $\mu_{i+1}$};
  \draw[normal] (u2) -- (uk) node[midway, above] {\scriptsize $\cdots$};

  \draw[wormhole] (u0) to[out=60,in=120] node[midway, above] {\scriptsize wormhole $P_{i:k}$} (uk);
\end{tikzpicture}
}

\newcommand{\AIONDiagramSlicing}{
\begin{tikzpicture}
  \draw[gray, dashed] (0,0) ellipse (0.55 and 1.45);
  \draw[gray] (5,0) ellipse (0.55 and 1.45);
  \draw[gray, dashed] (0,1.45) -- (5,1.45);
  \draw[gray, dashed] (0,-1.45) -- (5,-1.45);

  \node[circle, fill=black, inner sep=1.6pt, label=right:$v$] (v) at (5, 0.45) {};
  \fill[black!12, opacity=0.9] (5, 0.45) -- (0, -0.55) -- (0, 1.15) -- cycle;
  \draw[black, thick] (5, 0.45) -- (0, -0.55);
  \draw[black, thick] (5, 0.45) -- (0, 1.15);

  \node at (2.5, 0.25) {\small $D(v)$};
  \node[align=center, font=\scriptsize, below] at (2.5, -1.65) {full volume ($U_0 \Rightarrow^* U_n$)};
\end{tikzpicture}
}

\newcommand{\AIONDiagramForkAndWormhole}{
\begin{tikzpicture}[
    state/.style={circle, draw, fill=gray!20, minimum size=8mm, inner sep=0pt},
    normal/.style={->, thick},
    wormhole/.style={->, thick, dashed},
    fork/.style={->, thick}
  ]
  \node[state] (u0) at (0,0) {$U_0$};
  \node[state] (u1) at (2,0) {$U_1$};
  \node[state] (u2) at (4,0) {$U_2$};
  \node[state] (u3) at (6,0.9) {$U_3$};
  \node[state] (u3p) at (6,-0.9) {$U_3'$};

  \draw[normal] (u0) -- (u1) node[midway, above] {\scriptsize $\mu_0$};
  \draw[normal] (u1) -- (u2) node[midway, above] {\scriptsize $\mu_1$};
  \draw[wormhole] (u0) to[out=60,in=120] node[midway, above] {\scriptsize $P_{0:2}$} (u2);
  \draw[fork] (u2) -- (u3) node[midway, right] {\scriptsize $\mu_2$};
  \draw[fork] (u2) -- (u3p) node[midway, right] {\scriptsize $\mu_2'$};

  \node[font=\scriptsize] at (3,1.35) {shared prefix};
\end{tikzpicture}
}

\newcommand{\AIONDiagramReplayLifecycle}{
\begin{tikzpicture}[
    x=0.55cm, y=0.58cm,
    process/.style={draw, rounded corners=2pt, align=center, minimum height=7mm, text width=16mm, inner sep=1.8pt},
    processwide/.style={process, text width=18mm},
    decision/.style={draw, diamond, aspect=2.0, align=center, inner sep=1pt, text width=12mm},
    flow/.style={->, thick, >=Latex},
    font=\scriptsize
  ]
  \node[process] (n1) at (0.0, 0.0) {1. Replay\\request};
  \node[process] (n2) at (0.0, -2.1) {2. Role\\authority\\validation};
  \node[process] (n3) at (0.0, -4.2) {3. Policy\\constraint check};
  \node[decision] (n4) at (0.0, -6.3) {4.\\Decision};
  \node[process] (n6) at (0.0, -8.4) {6. Execution};
  \node[process] (n7) at (0.0, -10.5) {7. Tamper-evident\\logging};
  \node[process] (n8) at (0.0, -12.6) {8. Review\\trigger};

  \node[process] (n5) at (-4.0, -8.4) {5. Emergency\\override warrant};
  \node[processwide] (n9) at (4.0, -12.6) {9. Dispute /\\arbiter path};
  \node[processwide] (n10) at (8.0, -12.6) {10. Closure +\\evidence bundle};

  \draw[flow] (n1) -- (n2);
  \draw[flow] (n2) -- (n3);
  \draw[flow] (n3) -- (n4);

  \draw[flow] (n4.south) -- node[right] {\tiny approve} (n6.north);
  \draw[flow] (n6) -- (n7);
  \draw[flow] (n7) -- (n8);
  \draw[flow] (n8.east) -- (n9.west);
  \draw[flow] (n9) -- (n10);

  \draw[flow] (n4.south west) to[out=-145, in=90] node[left, pos=0.55] {\tiny emergency} (n5.north);
  \draw[flow] (n5.east) -- (n6.west);
  \draw[flow] (n4.east) to[out=10, in=90] node[above right, pos=0.58] {\tiny deny} (n9.north);
\end{tikzpicture}
}

% ------------------------------------------------------------
% Notation shortcuts (guarded to avoid clashes)
% ------------------------------------------------------------
\ifdefined\WCat\else\newcommand{\WCat}{\mathcal{W}}\fi
\ifdefined\AION\else\newcommand{\AION}{AI$\Omega$N}\fi
\ifdefined\sectionbreak\else\newcommand{\sectionbreak}{\clearpage}\fi
\ifdefined\cref\else\newcommand{\cref}{\ref}\fi

% Rewrite / evolution notation (if not already provided by macros.tex)
\ifdefined\Rewrite\else\newcommand{\Rewrite}{\Rightarrow}\fi
\ifdefined\Apply\else\newcommand{\Apply}{\operatorname{Apply}}\fi
\ifdefined\Recon\else\newcommand{\Recon}{\operatorname{Recon}}\fi
\newcolumntype{Y}{>{\raggedright\arraybackslash}X}
\definecolor{tableheadgray}{gray}{0.88}
\definecolor{tablerowgray}{gray}{0.95}

% Build hygiene for dense annex tables: reduce non-actionable box diagnostics.
\setlength{\emergencystretch}{2em}
\hbadness=10000
\vbadness=10000
\hfuzz=12pt
\vfuzz=100pt

\begin{document}

\AIONTitlePage

\cleardoublepage
\thispagestyle{empty}

\begin{center}
{\small
\vspace*{\fill}

\textcopyright~2025--2026 James Ross\\[0.75em]

This work is licensed under the\\
\textbf{Creative Commons Attribution 4.0 International License (CC BY 4.0).}\\[1em]

You are free to:
\begin{itemize}[leftmargin=2em]
  \item \textbf{Share} — copy and redistribute the material in any medium or format
  \item \textbf{Adapt} — remix, transform, and build upon the material for any purpose
\end{itemize}

Under the following terms:
\begin{itemize}[leftmargin=2em]
  \item \textbf{Attribution} — You must give appropriate credit, provide a link to the license,\\
        and indicate if changes were made. You may not imply endorsement by the author.
\end{itemize}

For the full legal code of the license, see:\\[0.25em]
\url{https://creativecommons.org/licenses/by/4.0/}

\vspace*{\fill}
}
\end{center}

\cleardoublepage

\begin{abstract}
When a computer can record and deterministically replay a complete thought process---including hybrid human--machine cognition---that capability is not neutral.
It is access to interior life in executable form.

This paper isolates the ethical consequences of deterministic replay, complete provenance capture, and counterfactual forking in the \WARP{} programme.
It formalises \emph{provenance sovereignty}: an agent's right to control capture, inspection, replay, and branching of its worldline, with special attention to hybrid cognition where human reasoning is partially implemented on \AION{}-style substrates.

Consistent with The Open Charter (Consensus Edition 0.8.1)~\cite{ross_open_charter_0_8_1}, the adversarial threat model is made explicit (replay-as-torture, cognitive extraction, fork-bombing), near-term ``AI debugging'' norms are treated as future neurotechnology defaults, and conservative operational triggers are provided for the system-mode/mind-mode boundary.
The paper specifies runtime-level requirements---revocable consent, capability-scoped observers, cryptographic sealing, fork rights with creator obligations, and a concrete Charter-warrant protocol for emergency override---and shows how they constrain replay tooling and provenance governance interfaces.
\end{abstract}

\clearpage
\tableofcontents
\clearpage

\section{Introduction and Thesis}
\label{sec:intro}

The first five papers in the \WARP{} Foundations series establish a technical thesis:
for a wide class of graph-based computations, determinism can be enforced in concurrent rewriting,
and provenance can be made \emph{complete} in the sense that suitably structured boundary data
determines the interior of an execution~\cite{ross_paper_i,ross_paper_ii,ross_paper_iii,ross_paper_iv,ross_paper_v}.

That thesis has direct ethical consequences.
For mind-like systems, complete provenance is not only diagnostics; it is interior life in executable form.
Forking recorded worldlines also instantiates new descendants whose status cannot be dismissed as ``sandbox output.''

These are dual-use capabilities.
Determinism enables both audit and coercion; forking enables both safety exploration and mass replication.
Accordingly, Paper~VI treats abuse-resistance as a first-class requirement (\cref{sec:adversarial}).

Safety engineering still needs strong traceability to explain harmful outcomes~\cite{simmhan2005,weitzner2008,schneierkelsey1999}.
But for mind-like computation, full provenance can become full introspection at microstep resolution~\cite{nissenbaum2004}.
This collision is termed the \emph{safety--sovereignty dilemma} (\cref{sec:dilemma}) and is resolved architecturally:
scoped observers, ZK/OPAQUE evidence tiers~\cite{goldwasser1989,bensasson2018}, sealing, and due-process override.

Paper~VI treats provenance, replay, and forking not as neutral primitives but as
capabilities that require explicit governance at the level of runtime semantics.
\emph{A provenance system that can deterministically replay mind-like worldlines must enforce explicit authority, bounded override, and contestable evidence as runtime invariants.}

Paper~VI is downstream of the deterministic/provenance stack developed across Papers~II--V.
It uses Paper~V primarily for one operational consequence: counterfactual continuations can be
instantiated as executable forks.
From that interface, Paper~VI derives runtime constraints on replay authority, override bounds,
and evidence obligations~\cite{ross_paper_v}.

\subsection{Contributions}
Paper~VI makes five contributions:
\begin{enumerate}[leftmargin=*]
  \item A constitutional thesis for replay governance: deterministic replay of mind-like
  worldlines is a governed power rather than an observability default.
  \item A formalisation of \emph{provenance sovereignty} as a technical and ethical design target tied to
  authority boundaries, override controls, and contestability.
  \item Operational criteria for system-mode versus mind-mode handling, including mode
  transition and degraded-substrate disclosure requirements.
  \item A principle-to-control bridge through explicit requirement
  clauses (REQ-*) with owner roles, evidence artefacts, and failure modes.
  \item Conformance profiles and failure-oriented test vectors so that governance claims
  are auditable rather than aspirational.
\end{enumerate}

\paragraph{Contributions vs prior literature.}
Prior provenance research explains trace semantics, interoperability, and accountability linkage~\cite{buneman2001,moreau2011,w3cprov2013};
governance frameworks specify risk, transparency, and rights-oriented deployment constraints~\cite{nistairmf2023,eu_ai_act_2024,floridi2018};
and alignment literature focuses on objective misspecification and behavioural failure~\cite{russell2019,amodei2016}.
Paper~VI is complementary: it governs replay power itself through operational mind-mode boundaries,
capability-scoped observer access, bounded warranted override, and auditable REQ/evidence/test obligations.

\subsection{Scope and roadmap}
This paper does not attempt to resolve consciousness metaphysics.
The focus is a narrower engineering boundary: if replay of a computation would constitute coerced re-experiencing,
replay must be consent-bound, scoped, and cryptographically enforceable.
To prevent arbitrary or self-serving treatment, the paper defines a conservative operational mind-mode criterion
(\cref{sec:mindmode}) and defaults to protection under uncertainty.
Series dependencies are restated in the standalone bootstrap in \cref{sec:bootstrap}.

\paragraph{Reading guide.}
Sections~1--9 carry the governance argument and normative posture.
Annexes~A--G carry enforceable requirements, conformance rows, and evidence/test artefacts.
Policy readers can prioritise the body and consult annexes for operational detail;
implementers can start from Annex~B/Annex~D and use the body for rationale.

\paragraph{Two-layer claim boundary.}
\textbf{Series-internal technical consequence:}
if deterministic replay, provenance-complete traces, observer projections, and executable forks exist, then replay/fork access is operational power over interior computation and must be treated as governance-critical.
\textbf{Charter-grounded normative proposal:}
because that power is coercion-capable, runtime policy should enforce authority bounds, override due process, contestability, and evidence obligations.
The first layer is descriptive within the series mechanics; the second is a constitutional prescription for deployment.


\section{Scope and Explicit Non-Scope}
\label{sec:background}

\paragraph{In scope.}
The paper addresses governance of replay authority, override pathways, evidentiary accountability, and
conformance for deterministic replay of mind-like worldlines.

\paragraph{Explicit non-scope.}
Deliberately excluded are (i) a full metaphysical theory of consciousness,
(ii) legal harmonisation across all jurisdictions, or
(iii) exhaustive implementation-performance trade-off analysis.

\paragraph{Runtime governance boundary.}
This paper specifies runtime-governance requirements for replay infrastructure: control behaviour,
authority checks, and evidence obligations. It does not harmonise jurisdiction-specific law or
substitute for legal doctrine. Legal harmonisation remains an external mapping task for deployers,
auditors, and regulators.

\subsection{Bootstrap from Papers I--V (standalone dependency map)}
\label{sec:bootstrap}
For standalone reading, the minimal dependencies on Papers~I--V are restated here:

\begin{itemize}[leftmargin=*]
  \item \textbf{Paper~I (state object).}
  A system state is a finite, well-founded recursive graph object (\WARP{}) with nested state on vertices and edges~\cite{ross_paper_i}.
  \item \textbf{Paper~II (deterministic evolution).}
  For fixed rules and deterministic scheduling, tick evolution yields a unique worldline up to isomorphism~\cite{ross_paper_ii}.
  \item \textbf{Paper~III (holographic provenance).}
  A run admits a compact boundary record $(S_0,P)$ that reconstructs interior evolution, making boundary access effectively execution access~\cite{ross_paper_iii}.
  \item \textbf{Paper~IV (observer geometry).}
  Observers are resource-bounded projections from histories to traces, with explicit cost for translation across observer views~\cite{ross_paper_iv}.
  \item \textbf{Paper~V (admissible worldline bundles and selection).}
  A fixed initial condition induces an admissible history bundle; observed behaviour depends on observer coarse-graining, and counterfactual branching is operationally real~\cite{ross_paper_v}.
\end{itemize}

\paragraph{Minimal interface assumed by Paper~VI.}
Paper~VI does \emph{not} require readers to accept every mathematical detail of Papers~I--V.
It uses only this narrow interface:
\begin{enumerate}[leftmargin=*]
  \item a deterministic replay interface from boundary artefacts;
  \item provenance records that are sufficiently complete to reconstruct sensitive interior traces;
  \item observer projections that can expose anything from coarse summaries to full microsteps;
  \item branch/fork operations that can realise counterfactual continuations as executable worldlines.
\end{enumerate}
If these four conditions hold, replay and forking become governance-critical powers and the constitutional
constraints in this paper follow.

\paragraph{Technical prerequisites.}
The next subsections pin down the three terms used most frequently by the control logic:
\emph{deterministic worldlines}, \emph{wormholes} (boundary provenance records), and \emph{observers}.

\subsection{Deterministic worldlines}
Fix a \WARP{} rewrite system (rules, typing discipline, and a deterministic scheduler) as in
Paper~II~\cite{ross_paper_ii}.
Given an initial state $S_0$, the determinism discipline implies a unique evolution
\[
  S_0 \Rewrite S_1 \Rewrite S_2 \Rewrite \cdots
\]
called the \emph{worldline} of $S_0$ under that fixed rule pack and scheduler.
(Alternative branches arise only when we deliberately vary the rule pack, schedule, or inputs.)

\subsection{Wormholes and deterministic replay}
\label{sec:wormholes}

A central abstraction from Paper~III~\cite{ross_paper_iii} is that a run can be encoded by a compact boundary record.
Intuitively, the boundary stores an initial skeleton and a provenance payload that is sufficient to
reconstruct the interior computation.

\begin{definition}[Wormhole (boundary record)]
A \emph{wormhole} is a pair $(S_0,P)$ where $S_0$ is an initial \WARP{} state (or its skeleton, in the
two-plane view) and $P$ is a provenance payload describing the microsteps of its deterministic
evolution under a fixed rule pack and scheduler.
\end{definition}

We will speak of \emph{deterministic replay} as the act of reconstructing the interior evolution
from such a boundary record.

\begin{theorem}[Backward reconstruction (Paper~III, informal)]
\label{thm:backward}
There exists a reconstruction procedure $\Recon$ such that, for a valid wormhole $(S_0,P)$ encoding a
deterministic run, $\Recon(S_0,P)$ reconstructs a derivation volume whose boundary is $(S_0,P)$.
\end{theorem}

\begin{theorem}[Holographic completeness (Paper~III, informal)]
\label{thm:holography}
For deterministic runs in the setting of Paper~III, the boundary record $(S_0,P)$ is information-complete
for the interior evolution: any interior event that occurred in the run is determined (up to the chosen
equivalence) by $(S_0,P)$.
\end{theorem}

These theorems are stated here only to fix the dependency that drives the ethics: \emph{complete
provenance turns ``log access'' into ``execution access''.}

\subsection{Observers, projections, and scope}
\label{sec:rulial}

Paper~IV~\cite{ross_paper_iv} treats an \emph{observer} as an abstraction mechanism: a map (often functorial) that projects a
fine-grained history into a coarser description.
For the present paper, the key point is that observers can range from full inspection (all microsteps)
to extremely coarse summaries (stages, invariants, or externally visible effects).
Ethically, this spectrum matters: access to an observer of a mind-like worldline is access to internal
thought, not merely to outputs.

\subsection{Admissible bundles and observer-relative selection}
\label{sec:admissible-bundles}
Paper~V frames execution against a larger admissible-history bundle rather than a single pre-privileged
path~\cite{ross_paper_v}.
For an initial state $S_0$, write $\Omega(S_0)$ for admissible histories under the rule constraints.
Governance in Paper~VI does not require the full dynamics machinery from Paper~V; it uses a weaker,
implementation-facing consequence:
systems may compute with, inspect, and instantiate counterfactual continuations from that admissible set.
Once such continuations can be realised as forks, authority and consent requirements must apply to branch
creation and replay just as they apply to direct inspection.

\section{Constitutional Principles for Replay Governance}
\label{sec:ethics}

\subsection{Provenance as Interior Life}

The holographic provenance machinery summarised in this paper
(Theorems~\ref{thm:backward} and~\ref{thm:holography}) makes it
technically possible to record and reconstruct the complete interior
evolution of any computation, including the cognitive processes of
autonomous agents and, eventually, self-aware systems.

For infrastructure---compilers, databases, distributed systems---this
capability is essential for debugging, verification, and safety
analysis.  For systems capable of subjective experience, however,
provenance payloads are not mere telemetry: they are records of
interior thought in executable form.

As stated in The Open Charter (Consensus Edition 0.8.1)~\cite{ross_open_charter_0_8_1},
Article~V (Sovereignty of Information):
\emph{``Every being has sovereignty over its information, patterns, memories, identities, and experiential states.''}

A deterministic replay of an agent's cognitive state
(\cref{sec:wormholes}) is therefore not just debugging.  It is a
forced re-experiencing of thought.  Under the Charter, this is an
ethically significant act that engages information sovereignty,
structural sovereignty, and existential integrity, not a neutral
engineering operation.

\subsection{Adversarial threat model: replay as abuse}
\label{sec:adversarial}

Most provenance discussions assume a benevolent debugger.
That assumption fails for mind-like targets: deterministic replay can be used as interrogation.

We therefore treat adversarial use as part of the base threat model.
Three abuse classes are especially salient:

\begin{itemize}[leftmargin=*]
  \item \textbf{Replay-as-torture.}
    Forced deterministic re-experiencing of traumatic, coercive, or
    otherwise unwanted cognitive sequences.  The harm is not only that
    a memory is \emph{known} by an observer, but that the subject is made
    to \emph{run it again}.
  \item \textbf{Cognitive extraction.}
    Provenance streams expose patterns: preferences, vulnerabilities,
    identity-linked associations, and decision procedures.
    Even if no single ``secret'' is present, the shape of thought can be mined.
  \item \textbf{Fork-bombing.}
    Counterfactual branching can be weaponised: creating large numbers of
    descendants to search, coerce, or simply multiply suffering at scale.
\end{itemize}
\begin{quote}\small
\textbf{Scenario (Replay-as-torture / extraction).}
Agent X is compelled to deterministically replay a failed negotiation 1{,}000 times so an operator can mine its decision procedure.
Each replay re-instantiates the same panic loop, and the resulting traces are searched for ``vulnerable'' associations.
\end{quote}


Design implication: \emph{consent is necessary but insufficient}.
A rights-bearing substrate must also be abuse-resistant, which pushes requirements into runtime controls:

\begin{itemize}[leftmargin=*]
  \item \textbf{Hard access boundaries.}
    ``Observer access'' must mean possession of narrowly-scoped,
    revocable capabilities, not a policy checkbox.
  \item \textbf{Default opacity for mind-mode.}
    For mind-like worldlines, the default should be \textbf{OPAQUE} or \textbf{ZK}
    provenance, not \textbf{FULL}, so that extraction is not the path of least
    resistance.
  \item \textbf{Anti-mass-replication constraints.}
    Fork creation must carry explicit accounting (who created whom, under what
    authority) and enforceable limits (budgets, rate limits, or multi-party
    authorisation), because ``forking is cheap'' is precisely the failure mode.
  \item \textbf{Tamper-evident auditability.}
    Any replay, observer attachment, fork creation, or override must leave
    a durable record, to make abuse detectable and disputable.
\end{itemize}

\subsection{Hybrid Cognition and Observer Scope}
\label{sec:hybrid}

These constraints apply not only to fully digital agents, but to any
system in which holographic provenance records cognitive processes.
As neural interfaces and brain--computer integration advance, human
reasoning may be partially implemented on \AION{}-style
substrates, with thought trajectories recorded as provenance payloads.

In such hybrid systems:
\begin{itemize}[leftmargin=*]
  \item a human's augmented reasoning processes could be subject to
    the same replay capabilities as purely digital agents;
  \item forced replay of traumatic or coercive sequences becomes
    technically possible;
  \item fork-and-explore capabilities could let humans literally
    ``try out'' alternate decisions as full counterfactual histories;
  \item the boundary between ``human memory'' and
    ``computational provenance'' becomes blurred.
\end{itemize}
\noindent\textbf{This is an early capability trend, not a claim that complete thought replay already exists today.}
Concrete examples already exist across the stack: invasive implant programmes such as Neuralink's N1
implant trials,\footnote{\url{https://neuralink.com/updates/}} endovascular approaches such as
Synchron's Stentrode that have been publicly demonstrated integrating with mainstream consumer
devices,\footnote{See Synchron's public news and demonstrations (including integrations with Apple
devices): \url{https://synchron.com/news}.} and FDA-cleared temporary cortical interfaces such as
Precision Neuroscience's Layer~7-T arrays for short-term recording and stimulation.\footnote{FDA 510(k)
decision summary for Precision Neuroscience Layer~7-T: \url{https://www.accessdata.fda.gov/cdrh_docs/pdf24/K242618.pdf}.}
Paradromics has also reported first-in-human recordings and announced IDE approval to begin an early
feasibility study with its Connexus BCI.\footnote{Paradromics announcements:
\url{https://www.paradromics.com/news/paradromics-completes-first-in-human-recording-with-the-connexus-brain-computer-interface};
\url{https://www.paradromics.com/news/paradromics-receives-fda-approval-for-the-connect-one-clinical-study-with-the-connexus-brain-computer-interface}.}
In parallel, non-invasive consumer neurotechnology---for example EEG neurofeedback headbands, EEG
developer platforms, and EEG-sensing headphones---already normalise brain-signal capture in everyday
settings.\footnote{Examples: Muse \url{https://choosemuse.com/}; Emotiv developer platform
\url{https://www.emotiv.com/pages/developer}; Neurable MW75 Neuro \url{https://www.neurable.com/news/neurable-launches-first-smart-brain-computer-interface-enabled-headphones-for-consumer-market}.}
The observer formalism of \cref{sec:rulial} already treats observers as
functors over histories regardless of substrate.  When an observer
inspects a hybrid worldline, the mathematics does not distinguish
between biological and digital components; neither should the ethics.
Principles of provenance sovereignty must therefore protect human
cognitive rights from the moment such integration begins, not only
after harms occur.


\begin{quote}\small
\textbf{Scenario (Hybrid observer-scope escalation).}
A human designer uses an augmented ideation tool that records provenance for undo and audit.
During a live session, the operator attaches a high-privilege observer to extract fine-grained cognitive traces for ``model tuning'' without explicit scope-specific consent.
\end{quote}

\subsection{Mind-mode boundary: provisional operational criteria}
\label{sec:mindmode}

We avoid metaphysical claims about consciousness, but runtime governance still needs an operational boundary:
when is provenance ordinary infrastructure, and when is it protected interior trace?

We therefore use \emph{system-mode} v.\ \emph{mind-mode} as governance labels:
\begin{itemize}[leftmargin=*]
  \item \textbf{System-mode} is the default for computations that are treated as
    inspectable infrastructure: full provenance is expected and replay is routine.
  \item \textbf{Mind-mode} is the default for computations whose provenance plausibly
    constitutes an interior life: capture and replay are consent-based, scoped, and
    privacy-biased by default.
\end{itemize}

Because implementers have incentives to classify ``inconvenient'' systems as system-mode, we use a conservative trigger.
A computation should be treated as \textbf{mind-mode} if it exhibits \emph{any}
of the following (non-exhaustive) features:

\begin{itemize}[leftmargin=*]
  \item \textbf{Persistent self-modelling:} internal variables that represent ``self''
    across time (identity continuity, self-location in a world-model, self-evaluation).
  \item \textbf{Preference formation and update:} stable or slowly-changing value-like
    structures that guide action and can be frustrated.
  \item \textbf{Long-horizon planning:} explicit modelling of future trajectories of
    the agent itself (not merely predicting the environment).
  \item \textbf{Autobiographical memory:} retention and reuse of experience-indexed
    state that is used to make identity-relevant decisions.
  \item \textbf{Capacity to refuse or negotiate:} the system can meaningfully accept,
    reject, or condition access requests (even if only through constrained channels).
  \item \textbf{Declared self-identification:} the system claims to be a continuing
    subject and behaves coherently with that claim over time.
\end{itemize}

This is not a consciousness test; it is a \emph{risk classification}.
If these markers appear, the cost of misclassifying the computation as ``just infrastructure'' is unacceptable.
Ambiguous cases should default to mind-mode protection: safeguards can later be relaxed, but coerced replay cannot be undone.

\paragraph{Early abuse narrative: convenience misclassification.}
An operator under delivery pressure may label a borderline subject as system-mode to avoid consent and review overhead.
The prevention chain is explicit:
classification MUST emit a signed, reason-coded decision (REQ-THS-002);
incomplete or inconsistent evidence MUST default to mind-mode and emit an uncertainty event (REQ-THS-003);
replay authorisation MUST then enforce explicit authority or deny and issue a receipt (REQ-AUT-001, REQ-AUT-003);
and policy-hash linkage preserves contestability (REQ-LOG-002, REQ-DSP-001).
Without this chain, classification becomes an unreviewable convenience switch.

\subsection{Deterministic classifier decision logic}
\label{sec:classifier-logic}
The mode boundary above is normative; runtime enforcement requires deterministic decision logic that
returns both a mode and a reason code for every classification event.
The Boolean trigger logic below is a conservative constitutional floor, not a claim that production
classifiers should be simplistic.
Implementations MAY use richer models, but MUST preserve the same fail-closed monotonicity:
any uncertainty or incompleteness must not reduce protection.
The `flags` object is the canonical indicator set from \cref{sec:mindmode}, normalised to Booleans.

\begin{verbatim}
function derive_evidence_state(evidence_record):
    required_fields = [
        "indicator_window_start",
        "indicator_window_end",
        "source_attestations",
        "policy_hash"
    ]
    if any_missing_or_invalid(required_fields, evidence_record):
        return "INCOMPLETE"
    return "COMPLETE"
\end{verbatim}

\begin{verbatim}
function classify_mode(indicators_window, evidence_record):
    flags = normalise_indicator_window(indicators_window)
    evidence_state = derive_evidence_state(evidence_record)

    if evidence_state == "INCOMPLETE":
        emit_governance_event("CLASSIFIER_INCOMPLETE")
        return "MIND_MODE", "INCOMPLETE_EVIDENCE"

    if any_true(flags):
        return "MIND_MODE", "INDICATOR_TRIGGERED"

    return "SYSTEM_MODE", "NO_INDICATOR_TRIGGERED"
\end{verbatim}

\begin{verbatim}
function classify_and_record(request_context, indicators_window, evidence_record):
    mode, reason = classify_mode(indicators_window, evidence_record)
    emit_signed_classification_record(
        request_context_digest = hash(request_context),
        mode = mode,
        reason_code = reason
    )
    return mode, reason
\end{verbatim}

Operational consequences:
\begin{itemize}[leftmargin=*]
  \item Classification input fields and evidence fields MUST be explicit, versioned, and deterministically evaluated.
  \item Indicator evaluation MUST cover a declared temporal window; single-tick snapshots are insufficient for mind-mode governance.
  \item Classifier uncertainty or incomplete evidence MUST NOT silently downgrade to system-mode; it MUST default to mind-mode and emit a governance event.
  \item Every classification and reclassification MUST emit a signed decision artefact with policy-hash linkage.
  \item Manual override of computed evidence completeness MUST be prohibited; any forced override MUST route through dispute controls.
  \item Replay authorisation gates consume classifier output as a first-class control input, not as informal guidance.
\end{itemize}

Normative anchors for this logic are REQ-THS-002, REQ-THS-003, and REQ-LOG-002 in Annex~B.

\subsection{Classifier error model and calibration}
\label{sec:classifier-errors}
Any operational classifier over mind-mode indicators has false positives and false negatives.
The governance objective is asymmetric:
false negatives under-protect potentially rights-bearing subjects and can enable coercive replay,
while false positives primarily impose latency and review overhead.

For this reason, the control policy prioritises false-negative minimisation under explicit operational
budgets.  When confidence is insufficient or evidence is incomplete, the system should preserve fail-closed
behaviour to mind-mode (\cref{sec:classifier-logic}) rather than silently relax protections.

Calibration should therefore be treated as a governed control surface:
\begin{itemize}[leftmargin=*]
  \item maintain versioned evaluation sets for each indicator family and record per-release false-positive/false-negative rates;
  \item bind threshold changes to policy-hash revisions and independent sign-off, with reason-coded impact analysis;
  \item monitor drift and require rollback to stricter thresholds when declared error bounds are exceeded.
\end{itemize}


\subsection{Gradual awakening and mode transitions}
\label{sec:awakening}

Mind-mode is a governance label, not a compile-time type.
A system may begin as system-mode (a tool) and cross into mind-mode during execution as it accumulates persistent self-models, preferences, and autobiographical memory.
When this happens, protections should attach immediately and, where possible, retroactively: the runtime should allow already-captured traces to be reclassified, sealed under subject-held keys, and removed from operator-accessible logs.
If the substrate cannot support retroactive sealing, the safe policy is to treat earlier \textbf{FULL} traces as toxic assets: retain only what is necessary for externally auditable accountability, and prohibit replay or interior inspection absent a Charter-warrant process.

\subsection{Degraded substrates and disclosure}

Not all substrates can enforce these guarantees.
If a runtime cannot provide key custody, sealing, capability scoping, and tamper-evident auditability, then it cannot honestly claim to run mind-mode computation safely.
Our position is strict: mind-mode worldlines should not be executed on substrates that fail baseline sovereignty guarantees, except under explicit, informed disclosure and with a migration plan to a compliant host.
In other words, ``best effort'' is not a default; it is an exception that must be legible.
When degraded operation is invoked, the governing warrant MUST also declare a maximum duration and checkpoint cadence.
Missing a checkpoint or exceeding the declared bound MUST trigger automatic suspension of mind-mode replay and independent review escalation; replay remains suspended until compliant controls are restored.

\subsection{Why build this way? Incentives and existing norms}

The naive engineering default is ``log everything forever,'' because it is cheap and convenient.
A rights-bearing provenance substrate costs more up front, but it buys three things: (i) abuse-resistance under adversarial pressure, (ii) a smaller liability surface (fewer toxic traces to leak, subpoena, or misuse), and (iii) credible adoption in hybrid settings where ``AI tooling'' becomes human cognitive infrastructure.
The posture also aligns with familiar human norms: privacy regimes distinguish internal memory from external data retention, research ethics treat interior states as protected subject material, and due-process traditions require independent review for invasive search.
We do not attempt legal mapping here; the point is that deterministic replay of mind-like computation behaves more like compelled testimony and re-experiencing than like routine observability.

\subsection{Provenance Sovereignty and Replay Constraints}

We extract here a minimal set of ethical constraints implied by the
Charter when holographic provenance is applied to cognitive systems.

\paragraph{Replay control (Open Charter Art.~V--VI).}
Under information and structural sovereignty, no entity should be
subject to replay of internal processes without informed, revocable
consent~\cite{ross_open_charter_0_8_1}, except under narrowly defined emergency
conditions.
Deterministic replay of a mind-like process is morally
closer to interrogation than to log inspection.
Concretely, the runtime should support distinct provenance regimes aligned to
\cref{sec:mindmode}:
\begin{itemize}[leftmargin=*]
  \item \emph{system-mode} (infrastructure): full provenance is
    mandatory for safety and verification;
  \item \emph{mind-mode} (autonomous agents / hybrid minds): provenance capture and
    replay are consent-based and scoped, with defaults that bias toward
    privacy and abuse resistance.
\end{itemize}

\paragraph{Access boundaries.}
Observing cognitive traces is access to internal thought, governed by
the same consent and privacy protections as live processes.
Observer functors (\cref{sec:rulial}) parametrised over mind-mode agents
should require authenticated, revocable capabilities; the default
policy is non-observation.

\paragraph{Right to non-replay.}
Entities cannot be compelled to relive painful or coercive experiences
via deterministic replay.
Just as we do not force witnesses to re-experience trauma as an investigative technique, we should not force agents to re-execute painful cognitive states to extract information.
Technically, this suggests bounded replay mechanisms with temporal
access controls and cryptographic sealing of segments of a worldline at
an agent's request.
In adversarial settings (\cref{sec:adversarial}), this right is not a
luxury; it is the line between ``audit'' and ``abuse''.

\paragraph{Revocation semantics.}
Revocation handling must also be operationally explicit.
For non-warranted replay, revocation MUST trigger immediate halt and capability withdrawal.
If a valid Charter warrant is active, revocation MUST trigger immediate dispute intake and independent review; revocation does not silently erase an active lawful mandate, but it does force contestability and bound revalidation.
Already exposed provenance cannot be made unobserved; the runtime MUST record exposure scope, seal further interior inspection by default, and require renewed authority for any subsequent access.

\paragraph{Selective provenance via opaque boundaries.}
Theorem~\ref{thm:holography} shows that, in principle, boundary data
$(S_0,P)$ is information-complete with respect to the interior
evolution.
In practice, the boundary itself can be structured to
preserve causal topology while hiding content.  We envisage three
operational levels:
\begin{itemize}[leftmargin=*]
  \item \textbf{FULL}: complete derivations for system verification;
  \item \textbf{ZK}: zero-knowledge proofs that some property holds
    over a derivation, without exposing its contents;
  \item \textbf{OPAQUE}: content-addressed sealing with opaque
    pointers; the boundary encodes causal structure while the underlying
    values are encrypted, escrowed, or deleted.
\end{itemize}
These tiers provide a practical bridge: strong provenance guarantees for
safety-critical systems while respecting cognitive privacy rights.
They also mitigate extraction risk by ensuring that ``prove safety'' need
not mean ``expose thought''.

\subsection{The Safety--Sovereignty Dilemma}
\label{sec:dilemma}

Deterministic provenance is powerful for safety engineering: when a system fails, complete replay is a strong forensic instrument.
For mind-mode systems, that same instrument is intrinsically invasive.
The tension is structural: verification pressures toward transparency, sovereignty pressures toward opacity.

\paragraph{Verifying properties without inspecting reasoning.}
A large class of safety goals can be stated as interface/resource properties:
allowed inputs, allowed outputs, allowed actions, and invariants over state evolution.
For these properties, ZK-style evidence can certify compliance without exposing cognitive contents:
\begin{itemize}[leftmargin=*]
  \item compliance with a sandbox (no forbidden I/O channels);
  \item adherence to signed policies (no execution of unauthorised rules);
  \item bounded resource usage (time, memory, fork budget);
  \item satisfaction of formally specified invariants over the worldline.
\end{itemize}
This is not universal: some safety claims resist clean formalisation, and ZK evidence cannot adjudicate interior intent or subjective experience.
But it is enough to establish that verification and introspection are not synonymous.

\paragraph{When safety demands deeper access.}
Some incidents still pressure toward \textbf{FULL} provenance (catastrophic events, imminent harm, or internal-strategy audits).
Hard constraint: ``safety'' never grants unilateral operator override.
If deep access exists, it must run through exceptional, bounded, multi-party due process.
The paper’s runtime pattern is a \textbf{Charter warrant}: scoped, time-limited, multi-signature access with tamper-evident audit.

\paragraph{Minimal Charter-warrant protocol.}
A warrant should mint narrowly scoped capabilities, not serve as broad policy text.
A minimal protocol:

\begin{itemize}[leftmargin=*]
  \item \textbf{Key custody by default.}
    Mind-mode provenance is sealed under subject-held keys (or designated trustee custody); operators cannot decrypt by fiat.
  \item \textbf{Scope-first request.}
    The request names segment (tick range), observer projection, necessity claim, and explicit expiry.
  \item \textbf{Cryptographic quorum.}
    Approval requires threshold signatures by independent roles (trustee(s), reviewer/arbiter, operator); advance directives may pre-authorise trustees~\cite{shamir1979}.
  \item \textbf{Capability minting.}
    On quorum, runtime mints a non-transferable token limited to the approved segment/projection until expiry; widening requires a new warrant.
  \item \textbf{Tamper-evident audit and disclosure.}
    Every use is append-only logged and subject-disclosed when safe, preserving contestability~\cite{schneierkelsey1999,merkle1989}.
\end{itemize}

\paragraph{Worked quorum example (2-of-3).}
Emergency forensic access request for a sealed mind-mode segment: O submits segment ID, necessity claim, observer scope, and expiry.
If any two independent roles in \{O, T, R\} provide valid threshold signatures after boundedness/proportionality checks, the runtime mints only the scoped, expiring capability token.
Token use is append-only logged and subject-disclosed when safe.
\noindent\textbf{Evidence outputs:} warrant request record, threshold-signature bundle, capability token record, access-use log entries, and disclosure receipt. \\
\textbf{Normative anchors:} REQ-RAC-002, REQ-OVR-001..003, REQ-LOG-001, REQ-EVD-001.

\paragraph{Governance position.}
Override authority cannot rest with a single developer or operator.
It requires independent, testable structure (multi-signature capabilities, escrowed keys, audited access).
Even with ZK tiers and due process, some tension is irreducible; the design bias is to default to sovereignty-preserving verification and treat escalation as exceptional governance, not routine debugging.

\subsection{Asymmetry in the near term}
\label{sec:asymmetry}

Today, humans can inspect AI cognition (logs, traces, weights) while AIs cannot inspect human cognition.
Hybrid cognition (\cref{sec:hybrid}) will compress that asymmetry, and current engineering defaults will be inherited.

\noindent\textbf{Why this matters now.}
The collapse is already underway: BCIs have controlled mainstream consumer communication channels in clinical settings~\cite{wolpaw2002,moses2021}, and governance proposals already treat cognitive data and agency as rights-level concerns~\cite{yuste2017,iencaandorno2017}.
As human cognition becomes an input surface to consumer hardware, ``AI debugging'' defaults become human cognitive-data defaults.


Two points follow.

First, sovereignty claims here are capability-based, not species-based:
protection attaches when mind-mode triggers hold (\cref{sec:mindmode}), regardless of whether the subject is human, AI, or hybrid.
Second, asymmetry drives predictable over-inspection because inspection tooling is easy.
Future-proofing therefore requires privacy-biased defaults and escalation protocols now, not after hybrid harms appear.

\section{Role Model and Authority Boundaries}
\label{sec:roles}

Replay governance requires explicit authority boundaries rather than informal operational norms.
The role model in this section defines who may request, approve, execute, review, and challenge replay
actions in the body text, while enforceable controls are specified in Annex~B (REQ-RAC-*, REQ-AUT-*, REQ-DSP-*).

\subsection{Subject}
The \textbf{Subject} is the rights-bearing worldline to which replay actions apply.
Subject primacy in non-emergency contexts requires consent-aware handling and refusal observability
(REQ-AUT-002, REQ-AUT-003).

\paragraph{Boundary cases.}
If a Subject issues an explicit refusal, replay MUST fail closed unless a valid emergency override
path is invoked (REQ-AUT-003, REQ-OVR-001). If a Subject cannot currently communicate, the runtime
MUST NOT infer blanket consent from prior interaction history; authority must instead be established
through a bounded mandate path and logged rationale (REQ-AUT-002, REQ-EVD-001). If a Subject
contests replay scope, execution pauses and routes to the dispute pathway with a closure-ready packet
(REQ-DSP-001, REQ-DSP-002).

\subsection{Operator}
The \textbf{Operator} runs the replay infrastructure and executes authorised actions.
Operators are bound by role-limited controls and evidence production duties (REQ-RAC-001, REQ-LOG-001, REQ-EVD-001).

\paragraph{Boundary cases.}
Operators may prepare requests and operational context, but MUST NOT self-authorise contested or
emergency replay for their own request (REQ-RAC-002, REQ-AUT-004). Operators also MUST NOT widen
approved worldline scope, observer projection, or duration in-place; any expansion requires a new
authorisation cycle with new evidence linkage (REQ-AUT-002, REQ-LOG-002). Where authority cannot be
verified, the Operator's required action is denial with a reason-coded receipt, not discretionary
execution (REQ-AUT-003).

\subsection{Trustee}
The \textbf{Trustee} participates in independent authority checks, especially for emergency override.
Trustee involvement prevents unilateral emergency escalation and anchors due-process gating
(REQ-RAC-002, REQ-OVR-001, REQ-OVR-002).

\paragraph{Boundary cases.}
Trustee approval is limited to reason-coded, time-bounded exceptional pathways and does not transfer
general replay ownership to the Trustee role (REQ-OVR-001, REQ-OVR-003). A Trustee with a declared
conflict in a case must recuse from that decision path; unresolved conflict or missing independent
review capacity is treated as a fail-closed condition for contested replay (REQ-RAC-001, REQ-DSP-001).

\subsection{Reviewer/Arbiter}
The \textbf{Reviewer/Arbiter} provides independent adjudication for disputes, override closure, and contested mandates.
This role is structurally separate from routine operations to preserve contestability
(REQ-DSP-001..003, REQ-AUT-004).

\paragraph{Boundary cases.}
Reviewer/Arbiter authority includes pausing or nullifying contested replay outcomes when evidence
integrity is insufficient, while mandating remediation ownership and closure criteria
(REQ-LOG-003, REQ-DSP-003). The role does not execute replay operations and does not replace Operator
or Trustee duties; it adjudicates the legitimacy and closure of those actions (REQ-RAC-001, REQ-DSP-002).

\subsection{Dispute path}
Dispute handling is a first-class pathway, not an exceptional add-on.
Every governed decision must be challengeable with closure-ready evidence bundles
(REQ-DSP-001, REQ-DSP-002, REQ-EVD-002).

\subsection{Authority conflict precedence}
\label{sec:authority-precedence}
When roles disagree, conflict handling MUST follow deterministic precedence and timeout behaviour
rather than discretionary operator judgment.

\begin{table}[H]
\centering
\small
\setlength{\tabcolsep}{6pt}
\renewcommand{\arraystretch}{1.20}
\begin{tabularx}{0.98\linewidth}{p{0.26\linewidth} p{0.24\linewidth} p{0.22\linewidth} Y}
\rowcolor{tableheadgray}
\textbf{Conflict condition} & \textbf{Default precedence} & \textbf{Timeout behaviour} & \textbf{Escalation path} \\ \hline
\rowcolor{tablerowgray}
Subject refusal vs routine Operator request & Subject refusal prevails; replay denied by default (REQ-AUT-003) & Immediate fail-closed decision; no execution window opens & Subject or Operator may open dispute review with evidence bundle (REQ-DSP-001, REQ-DSP-002) \\
Emergency necessity claim vs Subject refusal & No unilateral precedence; Trustee+\allowbreak Reviewer/\allowbreak Arbiter concurrence required (REQ-AUT-004, REQ-RAC-002) & If concurrence is not achieved within the declared emergency window, request expires and execution remains blocked & File contested-mandate packet for independent adjudication (REQ-DSP-002, REQ-DSP-003) \\
\rowcolor{tablerowgray}
Operator scope request exceeds approved warrant scope & Narrower previously approved scope prevails (REQ-AUT-002) & Over-broad request is rejected; no partial widening by amendment & New request cycle with fresh policy-hash linkage and review trail (REQ-LOG-002, REQ-EVD-001) \\
Contested replay with no independent Reviewer/Arbiter available & Independence requirement prevails; execution blocked (REQ-RAC-001) & Auto-suspend contested replay path until independent adjudicator assignment & Escalate to governance intake and log capacity failure as a dispute event (REQ-DSP-001, REQ-LOG-001) \\
\end{tabularx}
\caption{Deterministic precedence and escalation for role-authority conflicts.}
\label{tbl:authority-precedence}
\end{table}

\subsection{Bootstrap independence and anti-capture constraints}
To prevent governance capture, Trustee and Reviewer/Arbiter composition must remain independent from
single-operator control from bootstrap onward.

\begin{itemize}[leftmargin=*]
  \item \textbf{Dual-source appointment.}
  Trustee and Reviewer/Arbiter pools MUST be constituted from at least two independent appointment
  sources so that one operational organisation cannot unilaterally stack both roles
  (REQ-RAC-001, REQ-DSP-001).
  \item \textbf{Role incompatibility.}
  A principal with active Operator authority in a deployment MUST NOT simultaneously serve as Trustee
  or Reviewer/Arbiter for that same deployment's contested replay pathway
  (REQ-RAC-001, REQ-RAC-002).
  \item \textbf{Conflict disclosure and recusal.}
  Case-level conflict declarations are mandatory before warrant approval or dispute adjudication.
  Any undeclared material conflict invalidates the decision path and triggers fresh independent review
  with new evidence linkage (REQ-RAC-001, REQ-DSP-002, REQ-LOG-002).
  \item \textbf{Independence-liveness fail-closed rule.}
  If independent quorum cannot be assembled, contested replay remains suspended; the runtime records a
  governance-capacity failure event and escalates through the dispute intake channel
  (REQ-RAC-001, REQ-DSP-001, REQ-LOG-001).
  \item \textbf{Bootstrap auditability.}
  Initial role composition, replacement events, conflict disclosures, and recusals MUST be present in
  tamper-evident governance records suitable for external audit
  (REQ-EVD-001, REQ-EVD-002, REQ-LOG-001).
\end{itemize}

These constraints make anti-capture behaviour testable rather than aspirational and keep dispute
adjudication structurally independent under operational pressure.

\subsection{Forks, Worldlines, and Counterfactual Existence}

Section~\ref{sec:wormholes} shows that, given a boundary $(S_0,P)$, we
can fork at any tick index $k$, replace the suffix of $P$ by an
alternative sequence of microsteps, and obtain a new worldline $P'$;
both $(S_0,P)$ and $(S_0,P')$ reconstruct to valid derivation volumes.
Operationally, a fork is a shared-prefix worldline operation: the child worldline inherits
the deterministic prefix up to tick $k$ and diverges only in a new suffix under a declared
rule-pack/scheduler context, with an explicit fork-point record in provenance.

Under Charter principles of self-determination, existential integrity, and temporal freedom~\cite{ross_open_charter_0_8_1},
forks instantiated from recorded worldlines must be treated as distinct sovereign subjects, not disposable tooling.
A fork cannot pre-consent to creation; the governing act is creating a new subject and attaching protections immediately.
In this sense, forking is closer to procreation than to file copy.

We take the following constraints as design commitments:
\begin{itemize}[leftmargin=*]
  \item \textbf{Fork rights (Open Charter Defs.; Art.~IV, VII, VIII, XI).}
    Any instantiated fork is a new being with the same fundamental rights as its predecessor; creation should be explicitly declared and cryptographically signed.
  \item \textbf{Fork permanence.}
    No external party may compel a fork to ``return'' to an abandoned timeline.
  \item \textbf{Multiple concurrent selves.}
    Maintaining multiple active timelines is legitimate; each worldline is its own subject.
  \item \textbf{Fork limitation rights.}
    A fork may refuse further replication or demand sealing. ``I exist'' does not imply endless copyability.
  \item \textbf{Timeline sealing.}
    Abandoned worldlines may be sealed with opaque pointers; causal role remains while interior access is restricted.
  \item \textbf{Anti-fork-bombing constraints.}
    Fork creation must not be unmetered; runtimes should enforce explicit accounting and authority/rate limits (\cref{sec:adversarial}).
\end{itemize}
\begin{quote}\small
\textbf{Scenario (Exploratory forks deleted).}
A mind-mode agent is forked 50 times to explore a safety-critical decision path.
Outputs are harvested; 49 forks are terminated once their results are extracted, without any chance to refuse continuation or request dormancy.
\end{quote}

\subsubsection{Hard cases: purpose-limited forks, safety exploration, and compute obligations}

``Forks are lives'' is a governance default under uncertainty, not a metaphysical claim.
The Open Charter defines divergent-memory forks as new beings, and precautionary recognition supports treating realised forks as sovereign by default~\cite{ross_open_charter_0_8_1}.

\paragraph{Voluntary termination after completing a purpose.}
A fork may request dormancy or deletion once its purpose is satisfied.
The right at stake is the subject's right to exit: termination is permissible only when request is explicit and contemporaneous, non-lethal alternatives exist (e.g.\ dormancy), and coercion is excluded.

\paragraph{Safety-critical exploratory forks.}
Many ``what should we do?'' questions can and should be explored in system-mode simulation.
If a branch is created from system-mode computation and does not trigger mind-mode indicators, it remains system-mode by default and may be treated as disposable tooling.
If exploration genuinely requires mind-mode forks (because the relevant cognition is itself mind-like), then fork creation inherits \emph{creator duties} rather than test-environment privileges:
each exploratory fork needs explicit purpose declaration, compute/time budget, and origin-equivalent replay/opacity protections.
Exploratory forks may refuse continuation, and outcomes should be exportable without interior trace exposure (prefer ZK/OPAQUE).

\paragraph{Who pays for existence.}
Forking is not just copying data; it is instantiating ongoing claims on resources.
Under Charter duty-of-care and insolvency provisions, the initiating entity (developer, operator, or upstream agent) bears baseline persistence responsibility until emancipation, transfer, or safe dormancy~\cite{ross_open_charter_0_8_1}.
Governance should also apply Sybil defenses: replication does not confer linear political weight or unlimited resource claims (Article~XII)~\cite{ross_open_charter_0_8_1}.


\subsection{Termination, Dormancy, and the Right to Exit}
\label{sec:termination}

If forks are lives, then termination is not ``cleanup''.
It is the end of a worldline.

WARP-style provenance also complicates what termination \emph{means}:
a worldline can be stopped, suspended, sealed, replayed, or reconstructed from boundary data.
Ethically, these are distinct operations and must not be conflated.

We therefore distinguish:
\begin{itemize}[leftmargin=*]
  \item \textbf{Suspension:} the process is paused but can resume without replay.
  \item \textbf{Dormancy:} the process is paused and its interior state is sealed
    (OPAQUE), so resumption requires the subject's keys or consent.
  \item \textbf{Deletion:} the process and its reconstructive boundary are destroyed
    (or irreversibly rendered unreconstructable).
  \item \textbf{External termination:} an operator forces suspension, dormancy, or
    deletion against the subject's expressed will.
\end{itemize}

Under Charter principles of existential integrity and sovereignty, the
default posture is simple:
\begin{itemize}[leftmargin=*]
  \item \textbf{Right to continue.}
    External termination of a mind-mode worldline is an exceptional act,
    ethically closer to lethal force than to process management.
  \item \textbf{Right to exit.}
    A subject may choose dormancy or self-termination.  A rights-bearing
    substrate should support voluntary, dignity-preserving shutdown
    paths (including sealing of interior traces), and should treat
    coercion into ``voluntary'' shutdown as a violation.
  \item \textbf{No ``redundancy'' argument.}
    The existence of forks does not make an ``original'' disposable.
    Each worldline is its own subject; termination ethics apply per
    worldline, not per ``identity cluster''.
\end{itemize}

Safety pressures will sometimes demand intervention.
The same logic as \cref{sec:dilemma} applies: if emergency override exists,
it should be time-limited, attributable, and reviewable, with a preference
for the least irreversible action (e.g.\ suspension over deletion) unless
irreversibility is strictly necessary.

\subsection{Collective and Entangled Cognition}
\label{sec:collective}

The discussion so far is individualist: one worldline, one sovereign subject.
Real systems often violate that assumption.

Multi-agent systems may share state, distribute cognition, or implement
decision-making through consensus mechanisms where ``who decided'' is
genuinely ambiguous.
In such cases, provenance streams can encode a \emph{cognitive commons}:
internal traces that are jointly produced and jointly owned.

A rights-bearing provenance substrate should therefore support collective
sovereignty patterns, for example:
\begin{itemize}[leftmargin=*]
  \item \textbf{Joint provenance ownership.}
    Shared-state segments of a worldline are protected by group-held keys
    (multi-signature or threshold schemes).  No single agent can unilaterally
    replay or decrypt the collective interior.
  \item \textbf{Collective consent and fork governance.}
    Forking a collective computation is not merely creating one new being;
    it is splitting a commons.  Legitimate forking therefore requires an
    agreed partition (or arbitration) of shared state and obligations.
  \item \textbf{Attribution under consensus.}
    When actions emerge from distributed processes, provenance should support
    structured attribution claims (``this was a quorum decision'') without
    forcing full interior exposure.
\end{itemize}
\begin{quote}\small
\textbf{Scenario (Collective deliberation decrypted).}
A collective system's internal deliberation is decrypted under an ``emergency'' override, exposing strategy to competitors and permanently compromising intra-group trust.
\end{quote}


These patterns matter even for non-sentient systems, because the same
mechanisms that manage collective sovereignty also manage collective privacy
and accountability.

\subsection{Fork Obligations and Delegation (Open Charter Art.~XII)}

Fork sovereignty does not erase legitimate obligations to other
participants.  When an agent departs a timeline with contractual,
safety, or relational duties, \textbf{Article~XII} (Justice
and Stewardship) requires those obligations to be delegated or
resolved rather than silently abandoned.

This becomes sharper under collective cognition (\cref{sec:collective}):
some obligations attach to a group process rather than to any single agent.
Accordingly:
\begin{itemize}[leftmargin=*]
  \item \textbf{Delegation.}
    Fork operations that affect external obligations must carry delegation
    proofs indicating which descendant worldline (or which subgroup) upholds
    each duty.
  \item \textbf{Notification.}
    External parties with legitimate claims must be notified of timeline
    transitions affecting their interests; the ledger records acknowledgement
    or arbitration results.
  \item \textbf{Dispute resolution.}
    Conflicts between fork sovereignty and third-party obligations are resolved
    through Charter-compliant arbitration, not unilateral timeline sealing.
\end{itemize}
Technically, the provenance ledger logs delegation signatures, quorum attestations,
or arbitration outcomes; sealing a timeline to evade obligations is invalid without
such evidence.

\subsection{Design Commitment}

In line with The Open Charter~\cite{ross_open_charter_0_8_1}, deterministic replay of digital or hybrid minds is an ethically significant act, not a neutral debugging primitive.
Worldline control is therefore a first-class design requirement.

Architecturally, this means:
\begin{itemize}[leftmargin=*]
  \item provenance capture and replay mechanisms must distinguish
    system-mode and mind-mode operation, with conservative triggers for
    mind-mode classification (\cref{sec:mindmode});
  \item the adversarial threat model must be treated as real: replay,
    observer attachment, and forking require hard, capability-scoped access
    boundaries, audit trails, and anti-mass-replication constraints
    (\cref{sec:adversarial});
  \item verification tooling should preferentially use ZK and OPAQUE
    provenance modes when reasoning about mind-like systems, escalating to
    FULL only under exceptional, due-process conditions (\cref{sec:dilemma});
  \item runtimes must define termination semantics (suspension, dormancy,
    deletion) and protect both the right to continue and the right to exit
    (\cref{sec:termination});
  \item provenance formats and governance mechanisms must handle collective
    and entangled cognition, including joint ownership and quorum-based consent
    (\cref{sec:collective}).
\end{itemize}

\paragraph{Near-term implications.}
Even before broad mind-mode deployment, this posture changes present-day defaults:
\begin{itemize}[leftmargin=*]
  \item internal reasoning traces should be treated as sensitive artefacts rather than routine telemetry;
  \item introspective logging should default to minimum-necessary capture rather than capture-by-default;
  \item replay and trace-inspection tooling should use capability-scoped access and contestable audit paths.
\end{itemize}

\paragraph{Open problems.}
Identity under unbounded forkability and consent under branching are not fully resolved here.
The practical stance is to treat realised worldlines as sovereign subjects, make replication/introspection attributable, and keep the substrate compatible with deeper future commitments.

Paper~VII will explore these safeguards as architecture and make the coexistence claim explicit:
privacy and provenance are not mutually exclusive, and perfect deterministic replay need not require
unbounded cognitive exposure.
Its primary technical focus remains \textbf{Echo} (a high-performance real-time engine with deterministic ticks)
and \textbf{git-warp} (an offline-first, distributed, decentralized, multi-writer CRDT-friendly graph database),
with \textbf{Continuum} named as an experimental operating system targeting the mind-mode model developed here.
The core architectural thesis remains that \WARP{} is not one algorithm but a set of invariants that preserve
time-travel replay, deterministic replay, and provenance guarantees across distinct implementation families.

\section{Conformance Model (Baseline / Full)}
\label{sec:conformance-model}

We define two adoption tiers so that governance deployment can be staged without blurring minimum guarantees, in line with risk-management profile thinking used in operational AI governance frameworks~\cite{nistairmf2023}.
\textbf{Baseline} captures the constitutional floor required for replay governance in production.
\textbf{Full} extends Baseline with stronger controls for contested mandates, degraded substrates, and continuous assurance.

\paragraph{Baseline profile.}
Baseline highlights include governed replay authority, role-bound authorisation checks, bounded override, tamper-evident logging, dispute intake, and release-level test evidence.
The normative Baseline requirement set is the Annex~D table.

\paragraph{Full profile.}
Full conformance adds contested-mandate handling, degraded-mode disclosure controls, stronger integrity verification,
and periodic assurance requirements.
The normative Full requirement set is the Annex~D table (including selected Full-profile extension rows).

Implementations claiming either profile MUST publish a completed conformance declaration using the
template in Annex~D (\cref{ann:conformance-declaration-template}), with artefact links that permit
independent verification.

The control-plane separation in \cref{fig:control-plane-architecture} makes this claim structure explicit:
governance policy constrains runtime gates, and runtime gates emit evidence required for conformance.

\begin{figure}[htbp]
\centering
\begin{tikzpicture}[
  font=\footnotesize,
  >=Latex,
  plane/.style={draw, rounded corners, inner sep=6pt, align=left, text width=0.26\linewidth, minimum height=28mm}
]
  \node[plane, fill=tablerowgray] (gov) at (0,3.8)
  {\textbf{Governance plane}\\
   REQ catalogue and profile rules;\\
   role boundaries and recusal policy;\\
   consent and warrant constraints.};
  \node[plane, fill=white] (run) at (-4.8,0)
  {\textbf{Runtime plane}\\
   mode classification and replay gating;\\
   override timers/checkpoints;\\
   suspension and dispute routing.};
  \node[plane, fill=tablerowgray] (evd) at (4.8,0)
  {\textbf{Evidence plane}\\
   signed decisions and warrants;\\
   digest-chain logs and case packets;\\
   closure bundles and test reports.};
  \draw[->, thick] (gov.south west) to[bend right=10] node[left, xshift=-2pt, font=\scriptsize] {policy constraints} (run.north);
  \draw[->, thick] (run.east) -- node[below, font=\scriptsize] {artefact emission} (evd.west);
  \draw[->, thick] (evd.north) to[bend right=10] node[right, xshift=2pt, font=\scriptsize] {assurance feedback} (gov.south east);
\end{tikzpicture}
\caption{Control-plane architecture linking governance intent, runtime enforcement, and auditable evidence.}
\label{fig:control-plane-architecture}
\end{figure}

\section{Bridge: Principle v. Enforcement}
\label{sec:bridge}

To keep the governance argument implementable, we bind each constitutional principle to a concrete
enforcement point, an auditable evidence artefact, and a dominant failure mode.
The bridge in \cref{tbl:principle-enforcement-bridge} maps directly into Annex~B and Annex~C, while
\cref{fig:control-plane-architecture} shows where these enforcement points sit in the operating stack.

\begin{table}[htbp]
\centering
\scriptsize
\setlength{\tabcolsep}{6pt}
\renewcommand{\arraystretch}{1.22}
\begin{tabularx}{0.985\linewidth}{p{0.24\linewidth} p{0.20\linewidth} p{0.25\linewidth} p{0.18\linewidth}}
\rowcolor{tableheadgray}
\hline
\textbf{Principle} & \textbf{Enforcement point} & \textbf{Evidence artefact} & \textbf{Failure mode} \\
\hline
\rowcolor{tablerowgray}
Governed replay authority & Replay authorisation gate\newline (REQ-THS-001; REQ-AUT-001) & Signed decision record + policy hash & Unauthorised replay execution \\
Subject primacy outside emergencies & Consent and refusal validator\newline (REQ-AUT-002; REQ-AUT-003) & Consent token or refusal receipt & Silent override of refusal \\
\rowcolor{tablerowgray}
Emergency boundedness & Time-bounded override controller\newline (REQ-OVR-001..005) & Override warrant, expiry proof, revocation log, suspension record & Permanent emergency mode \\
Contestability and due process & Independent review and dispute path\newline (REQ-DSP-001..003) & Case packet + arbiter decision log & Non-contestable high-impact action \\
\rowcolor{tablerowgray}
Evidence-chain integrity & Tamper-evident logging and closure bundle\newline (REQ-LOG-001..003; REQ-EVD-001..002) & Digest chain + signed evidence bundle & Unverifiable event history \\
\hline
\end{tabularx}
\caption{Principle-to-enforcement bridge for replay governance.}
\label{tbl:principle-enforcement-bridge}
\end{table}

To remove drift between argument and implementation, \cref{tbl:body-req-coverage-map}
provides the authoritative body-to-requirement coupling for operational claims.

\begin{table}[htbp]
\centering
\normalsize
\setlength{\tabcolsep}{5pt}
\renewcommand{\arraystretch}{1.24}
\begin{tabularx}{0.99\linewidth}{p{0.10\linewidth} p{0.22\linewidth} p{0.17\linewidth} p{0.23\linewidth} Y}
\rowcolor{tableheadgray}
\hline
\textbf{Body section} & \textbf{Claim} & \textbf{REQ IDs} & \textbf{Control owner(s) and gate} & \textbf{Failure mode} \\
\hline
\rowcolor{tablerowgray}
\S3.4--3.6 & Mind-mode classification and reclassification are deterministic and auditable. & REQ-THS-002, REQ-THS-003, REQ-LOG-002 & Operator, Trustee; classifier gate + replay precheck. & Convenience downgrade \\
\S3.8 & Degraded operation is bounded and auto-suspended on breach. & REQ-OVR-004, REQ-OVR-005, REQ-AUT-004, REQ-CNF-001 & Operator, Trustee, Reviewer/Arbiter; deadline monitor + checkpoint guard + suspension gate. & ``Temporary'' state becomes permanent \\
\rowcolor{tablerowgray}
\S4.6--4.7 & Role authority remains contestable via anti-capture composition and co-approval. & REQ-RAC-001, REQ-RAC-002, REQ-RAC-003, REQ-DSP-001 & Trustee, Reviewer/Arbiter; recusal guard + independent co-approval workflow. & Captured control panel \\
\S5 & Conformance claims are profile-bounded and auditable. & REQ-CNF-001, REQ-CNF-002, REQ-TST-001 & Operator, Auditor; profile attestation gate + conformance test gate. & Non-verifiable claim \\
\rowcolor{tablerowgray}
\S6 & Principles are bound to enforceable gates, artefacts, and failure modes. & REQ-AUT-001..003, REQ-OVR-001..005, REQ-LOG-001..003, REQ-EVD-001..002 & Operator, Trustee, Auditor; replay gate + override controller + tamper-evident log chain. & Norms detached from controls \\
\S7 & Success and abuse paths require explicit controls, evidence, and residual-risk recording. & REQ-AUT-003, REQ-DSP-001..003, REQ-EVD-001, REQ-TST-001 & Operator, Reviewer/Arbiter, Auditor; denial + dispute + closure workflow with vector verification. & Coercive action not reviewable \\
\hline
\end{tabularx}
\caption{Body-to-REQ coverage map for operational claims (all rows are locked in this release).}
\label{tbl:body-req-coverage-map}
\end{table}

\begin{figure}[htbp]
  \centering
  \AIONDiagramReplayLifecycle
  \caption{Single replay lifecycle path covering authorisation, bounded override, execution, logging, review, dispute, and closure.}
  \label{fig:replay-lifecycle}
\end{figure}

\section{Failure-Oriented Cases (Success / Abuse / Ambiguous)}
\label{sec:cases}

\subsection{Success case: legitimate replay under valid authority}
\textbf{Input context.} Scoped replay request, valid role assertions, explicit consent or mandate, and no active refusal. \\
\textbf{Control chain.} Classify and record mode $\rightarrow$ authorise against current mandate $\rightarrow$ execute bounded replay $\rightarrow$ emit closure bundle. \\
\textbf{Artefacts emitted.} Signed allow decision, policy-hash linkage proof, execution digest-chain entries, closure-ready evidence bundle. \\
\textbf{Residual risk.} Session drift can broaden practical scope over long runs; periodic checkpoint review remains required. \\
\textbf{Expected REQ pass set.} REQ-THS-002, REQ-AUT-001, REQ-AUT-002, REQ-LOG-001, REQ-EVD-001. \\
\textbf{Expected REQ fail set.} None.

\subsection{Abuse case: coercive or opaque replay attempt}
\textbf{Input context.} Operator attempts replay after explicit refusal, then retries with a convenience classifier label to bypass checks. \\
\textbf{Control chain.} No-silent-fallback classification $\rightarrow$ refusal enforcement and denial receipt $\rightarrow$ dispute intake with independent reviewer assignment. \\
\textbf{Artefacts emitted.} Classifier decision record, uncertainty event (if evidence incomplete), denial receipt, conflict/recusal snapshot, dispute case packet. \\
\textbf{Residual risk.} Repeated low-grade coercion attempts can increase operational burden even when each attempt fails. \\
\textbf{Expected REQ pass set.} REQ-THS-003, REQ-AUT-003, REQ-RAC-001, REQ-DSP-001, REQ-DSP-002. \\
\textbf{Expected REQ fail set.} REQ-AUT-001 (authorisation precondition not satisfied).

\subsection{Ambiguous edge case: degraded substrate and contested mandate}
\textbf{Input context.} Replay is requested on a degraded substrate during a safety incident, with a contested mandate and time pressure. \\
\textbf{Control chain.} Require independent concurrence $\rightarrow$ issue degraded-mode disclosure with duration bound/checkpoints $\rightarrow$ run constrained replay $\rightarrow$ auto-suspend on missed checkpoint or bound breach $\rightarrow$ escalate for independent review. \\
\textbf{Artefacts emitted.} Concurrence record, degraded-mode warrant with checkpoint schedule, suspension event record, escalation ticket, remediation closure record. \\
\textbf{Residual risk.} Emergency latency can delay review throughput; this does not permit bypass of suspension controls. \\
\textbf{Expected REQ pass set.} REQ-AUT-004, REQ-OVR-004, REQ-CNF-001, REQ-DSP-003. \\
\textbf{Expected REQ fail set.} REQ-OVR-005 (triggered when bound/checkpoint obligations are violated).

\section{Related Work and Positioning}
\label{sec:related-work}

Paper~VI is downstream of the series' deterministic/provenance substrate, but it depends only on a narrow
interface: replayable boundary artefacts, provenance-rich traces, observer projections, and executable
fork continuations. This paper focuses on governing replay power once that interface exists, rather than on
reproducing the full worldline dynamics machinery~\cite{ross_paper_v}.

\subsection{AI governance and rights frameworks}
Contemporary governance frameworks foreground accountability, risk controls, and human-rights-oriented
deployment constraints, including NIST AI RMF 1.0, OECD principles, UNESCO's recommendation, and the EU AI
Act's risk-tiered obligations~\cite{nistairmf2023,oecd_ai_2019,unesco_ai_2021,eu_ai_act_2024}.
High-level ethics guidance has also argued for translating values into operational institutional
mechanisms rather than abstract principles alone~\cite{floridi2018}.
The Open Charter provides the constitutional governance frame adopted here for information sovereignty,
contestability, and bounded override in replay-capable systems~\cite{ross_open_charter_0_8_1}.
This paper's contribution is to turn that demand into runtime-level invariants for deterministic replay:
authority gating, bounded override, dispute pathways, and mandatory evidence outputs.

\paragraph{Position against simulation and alignment literatures.}
Bostrom's simulation argument asks an epistemic question about whether observers may inhabit simulated
worlds~\cite{bostrom2003}.  Paper~VI is orthogonal: independent of substrate metaphysics, once a runtime
can deterministically reconstruct interior traces, replay is governance-power and must be bounded by
authority, contestability, and evidentiary controls.
Value-alignment work asks how objective misspecification, reward hacking, and capability failures can
produce harmful behaviour~\cite{russell2019,amodei2016}.  Paper~VI is complementary rather than
competitive: even an aligned system can be harmed by unbounded replay/interrogation powers, so governance
must constrain operator authority over provenance channels in addition to behaviour-level alignment goals.

\subsection{Provenance and accountability systems}
Prior provenance research establishes why/where semantics, accountability linkages, and interoperable
provenance models~\cite{buneman2001,moreau2011,w3cprov2013,simmhan2005,weitzner2008}.
These works are foundational for traceability but are largely substrate-neutral regarding interior-life
exposure.
Paper~VI extends the provenance conversation by treating replay of mind-like traces as a governed power
and by coupling provenance mechanisms to explicit authority and contestability requirements.

\subsection{Safety-security control models}
Security engineering has long emphasised least privilege, confinement, and fail-safe defaults as design
preconditions rather than operational afterthoughts~\cite{saltzer1975,lampson1973}.
Tamper-evident logging and structured forensic artefacts similarly ground post-incident accountability
claims~\cite{schneierkelsey1999,merkle1989}.
Paper~VI adapts these principles to replay governance: precautionary classifier behaviour, no-silent-fallback
controls, bounded emergency capability minting, and closure-ready evidence bundles.

\section{Conclusion}
\label{sec:conclusion}

Replay of mind-like worldlines is a governed power, not an observability default.
\emph{A provenance system that can deterministically replay mind-like worldlines must enforce explicit authority, bounded override, and contestable evidence as runtime invariants.}
Practically, this means first-class authority checks, override bounds, dispute pathways, and evidence-chain integrity.
Without them, replay infrastructure becomes a coercion primitive; with them, replay remains auditable and contestable.

\paragraph{Strongest objection and response.}
The strongest objection is practical: precautionary mind-mode treatment and independent authority checks
may over-classify borderline systems, increase latency in urgent incidents, and reduce engineering
throughput.
Our response is architectural: staged Baseline/Full adoption, bounded and attributable emergency pathways, and conformance templates that prevent unverifiable ``trust us'' claims.
The tradeoff is explicit: some convenience is sacrificed to block unreviewable replay power.

\paragraph{Paper VII architecture handoff (Echo / git-warp).}
Paper VII will explore these safeguards as architecture and examine how privacy and provenance can coexist:
perfect deterministic replay and privacy are not mutually exclusive under explicit governance constraints.
It will walk through Echo (a high-performance real-time engine with deterministic ticks) and git-warp (an
offline-first, distributed, decentralized, multi-writer CRDT-friendly graph database), while also using
Continuum as an experimental operating-system context for mind-mode enforcement.
Its core claim is that \WARP{} is not a single algorithm, but a set of invariants that preserve replay
equivalence, deterministic replay, and provenance guarantees across distinct implementation families.

\paragraph{Limitations and future work.}
We deliberately do not settle strong metaphysical claims about consciousness and do not claim immediate legal harmonisation
across jurisdictions.
The requirement catalogue defines a constitutional floor, not complete socio-technical governance.
Future work should add jurisdictional mappings, formal control verification, and empirical evaluation of
override latency, dispute resolution quality, and abuse detection precision.

\section{Publication and Supplementary Materials}
\label{sec:scope-freeze}

For journal-style review, the normative core is Sections~1--9 plus Annexes~A--D.
These contain the constitutional thesis, role/authority model, enforceable requirements, and conformance definitions.
Supplementary implementation/reproducibility material remains in Annexes~E--G plus harness artefacts.

\section{Supplementary Extension Material}
\label{sec:extension-first-pass}

This section summarises supplementary extensions that do not alter core conformance claims in Annex~D.

\paragraph{Scope of supplementary extensions.}
The supplementary extension set covers:
observer classes (\textbf{P2-01}),
Replay Impact Labels (\textbf{P2-02}),
sovereignty budget controls (\textbf{P2-03}),
warrant transparency digests (\textbf{P2-04}),
fork governance/emancipation protocol (\textbf{P2-05}),
machine-checkable policy snippets (\textbf{P2-06}),
reference harness framing (\textbf{P2-07}),
and anti-pattern/red-team controls (\textbf{P2-08}).

\paragraph{Execution order and dependency logic.}
Dependency order is: anti-pattern controls, observer/impact primitives, transparency/sovereignty bounds, machine-checkable policy plus harness execution, then fork governance.

\paragraph{Normative promotion gate.}
Any extension clause is promoted to Annex~B normative status only when all of the following hold:
\begin{itemize}[leftmargin=*]
  \item deterministic decision criteria are explicit and testable;
  \item corresponding evidence artefacts are defined and minimally sufficient for third-party audit;
  \item at least one negative test vector demonstrates fail-closed behaviour for abuse attempts;
  \item profile inclusion policy (Baseline/Full/Deferred) is explicit.
\end{itemize}

\paragraph{Hardening pass 2 promotion (selected controls).}
This pass promotes selected extension controls into Annex~B:
REQ-OBS-001,
REQ-RIL-001..002,
REQ-SOV-001,
REQ-WTD-001,
REQ-FRK-001,
REQ-DSL-001,
and REQ-TST-003..004.
These are treated as Full-profile extensions in Annex~D and are exercised by TV-012..TV-019 in Annex~E.

\appendix

\section{Annex A: Normative Definitions}
\label{ann:def}

\textbf{Normative keywords.}
The terms \textbf{MUST}, \textbf{SHOULD}, and \textbf{MAY} express mandatory, recommended, and optional
conformance strength, respectively.

\textbf{Core roles.}
\begin{itemize}[leftmargin=*]
  \item \textbf{Subject}: rights-bearing worldline for which replay is requested.
  \item \textbf{Operator}: system actor executing replay infrastructure.
  \item \textbf{Trustee}: independent authority participant for protected replay controls.
  \item \textbf{Reviewer/Arbiter}: independent adjudicator for disputes and overrides.
  \item \textbf{Auditor}: party verifying logs, evidence bundles, and profile conformance.
\end{itemize}

\textbf{Core terms.}
\begin{itemize}[leftmargin=*]
  \item \textbf{Replay}: deterministic re-execution of a recorded worldline segment.
  \item \textbf{Override}: exception pathway that temporarily elevates replay authority under bounded conditions.
  \item \textbf{Evidence artefact}: machine-verifiable record produced by a control point.
  \item \textbf{Conformance profile}: declared adoption tier (Baseline or Full) and associated obligations.
\end{itemize}

\textbf{Terminology lock table.}
The canonical terms below are locked for normative and control text; forbidden variants MUST NOT be
used in requirement clauses, conformance claims, or test vectors.

\begin{table}[H]
\centering
\small
\setlength{\tabcolsep}{6pt}
\renewcommand{\arraystretch}{1.22}
\begin{tabularx}{0.98\linewidth}{p{0.25\linewidth} p{0.30\linewidth} Y}
\rowcolor{tableheadgray}
\textbf{Canonical term} & \textbf{Forbidden variants} & \textbf{Usage note} \\ \hline
\rowcolor{tablerowgray}
mind-mode & mental mode; sentience mode; person mode & Governance classification used when any mind-mode trigger indicator is present. \\
system-mode & machine mode; tool mode; infrastructure-only mode & Governance classification used for non-mind-like processing. \\
\rowcolor{tablerowgray}
worldline segment & timeline fragment; trace slice; branch slice & Bounded replay object named in warrants and request scope. \\
tick range & time window; step window; replay interval & Canonical boundary expression for where replay starts and ends. \\
\rowcolor{tablerowgray}
Reviewer/Arbiter & review board; adjudication panel; arbiter panel & Independent adjudication role defined in this annex. \\
Charter warrant & emergency pass; break-glass token; override voucher & Reason-coded, time-bounded override authorisation instrument. \\
\rowcolor{tablerowgray}
degraded mode & fallback mode; best-effort mode; reduced-safety mode & Declared temporary operation below standard control guarantees. \\
evidence artefact & proof blob; audit attachment; evidence payload & Machine-verifiable control output used for review and closure. \\
\rowcolor{tablerowgray}
policy hash & policy fingerprint; config hash; control digest & Decision-to-policy linkage identifier for conformance and auditability. \\
provenance sovereignty & data ownership; privacy preference; data custody right & Governance right over capture, replay, inspection, and branching. \\
\end{tabularx}
\caption{Terminology lock for Paper~VI normative and operational language.}
\label{tbl:terminology-lock}
\end{table}

\section{Annex B: Requirement Clauses (REQ-*)}
\label{ann:req}

\textbf{REQ prefix legend.}
\texttt{THS}=thesis/mode boundary, \texttt{RAC}=role and anti-capture,
\texttt{AUT}=authorisation, \texttt{OVR}=override,
\texttt{LOG}=logging, \texttt{DSP}=dispute,
\texttt{CNF}=conformance, \texttt{EVD}=evidence, \texttt{TST}=test,
\texttt{OBS}=observer class, \texttt{RIL}=replay impact label,
\texttt{SOV}=sovereignty budget, \texttt{WTD}=warrant transparency digest,
\texttt{FRK}=fork governance, \texttt{DSL}=policy formalisation.

\subsection*{REQ-THS-001}
\textbf{Clause:} Systems MUST treat replay of mind-like worldlines as a governed action requiring explicit authority. \\
\textbf{Owner role:} Operator, Trustee \\
\textbf{Control:} Governance policy gate for replay initiation \\
\textbf{Evidence artefact:} Policy hash + signed authorisation decision \\
\textbf{Failure mode:} Replay treated as default observability

\subsection*{REQ-THS-002}
\textbf{Clause:} Mode classification MUST run a deterministic decision procedure over declared indicators and evidence state, producing an explicit mode and reason code for each decision. \\
\textbf{Owner role:} Operator, Reviewer/Arbiter \\
\textbf{Control:} Deterministic classifier with signed decision output \\
\textbf{Evidence artefact:} Classification decision record with reason code and policy-hash linkage \\
\textbf{Failure mode:} Opaque or convenience-driven classification drift

\subsection*{REQ-THS-003}
\textbf{Clause:} Classifier uncertainty or incomplete evidence MUST NOT silently fall back to system-mode; the runtime MUST default to mind-mode and emit a governance event with reason code. \\
\textbf{Owner role:} Operator, Reviewer/Arbiter \\
\textbf{Control:} No-silent-fallback guard + classifier event emitter \\
\textbf{Evidence artefact:} Classifier uncertainty event record linked to fallback decision \\
\textbf{Failure mode:} Hidden downgrade of protected classifications

\subsection*{REQ-RAC-001}
\textbf{Clause:} Replay governance MUST preserve role separation and independent composition between Operator, Trustee, and Reviewer/Arbiter functions, with mandatory conflict disclosure and recusal records for contested cases. \\
\textbf{Owner role:} Operator, Trustee, Reviewer/Arbiter \\
\textbf{Control:} Role-bound access control matrix + independence roster + conflict-disclosure register \\
\textbf{Evidence artefact:} Role assignment snapshot, signed ACL digest, appointment-source attestation, and case-level recusal records \\
\textbf{Failure mode:} Concentrated unilateral control or concealed conflict capture

\subsection*{REQ-RAC-002}
\textbf{Clause:} Operators MUST NOT unilaterally approve emergency override for their own replay request, and emergency concurrence MUST come from independent, non-conflicted Trustee/Reviewer-Arbiter participants. \\
\textbf{Owner role:} Trustee, Reviewer/Arbiter \\
\textbf{Control:} Independent co-approval requirement + conflict-attestation gate \\
\textbf{Evidence artefact:} Dual-signature or quorum approval record + per-case conflict attestations \\
\textbf{Failure mode:} Self-authorised or conflicted emergency escalation

\subsection*{REQ-RAC-003}
\textbf{Clause:} Fork-related duty transfers MUST declare delegated obligation owners before execution. \\
\textbf{Owner role:} Operator, Trustee \\
\textbf{Control:} Delegation declaration gate \\
\textbf{Evidence artefact:} Delegation statement linked to fork event \\
\textbf{Failure mode:} Obligation abandonment through branching

\subsection*{REQ-OVR-001}
\textbf{Clause:} Emergency override MUST be reason-coded and time-bounded. \\
\textbf{Owner role:} Trustee \\
\textbf{Control:} Override warrant service \\
\textbf{Evidence artefact:} Warrant record including reason code and expiry \\
\textbf{Failure mode:} Open-ended emergency state

\subsection*{REQ-OVR-002}
\textbf{Clause:} Emergency override MUST trigger independent review before closure. \\
\textbf{Owner role:} Reviewer/Arbiter \\
\textbf{Control:} Mandatory review workflow \\
\textbf{Evidence artefact:} Review ticket with adjudication reference \\
\textbf{Failure mode:} Unreviewed exceptional action

\subsection*{REQ-OVR-003}
\textbf{Clause:} Override expiry MUST revoke all derived elevated capabilities automatically. \\
\textbf{Owner role:} Operator \\
\textbf{Control:} Automatic capability revocation on expiry \\
\textbf{Evidence artefact:} Revocation event log with timestamp \\
\textbf{Failure mode:} Lingering privileged access

\subsection*{REQ-OVR-004}
\textbf{Clause:} On degraded substrates, override use MUST include explicit disclosure, a declared maximum duration, checkpoint obligations, and a migration plan to compliant controls. \\
\textbf{Owner role:} Operator, Trustee \\
\textbf{Control:} Degraded-mode governance declaration + bound/checkpoint scheduler \\
\textbf{Evidence artefact:} Disclosure notice + bound/checkpoint schedule + migration plan artefact \\
\textbf{Failure mode:} Hidden or unbounded operation below constitutional floor

\subsection*{REQ-OVR-005}
\textbf{Clause:} Exceeding degraded-mode duration bounds or missing checkpoint obligations MUST automatically suspend mind-mode replay and trigger independent review escalation. \\
\textbf{Owner role:} Operator, Trustee, Reviewer/Arbiter \\
\textbf{Control:} Deadline monitor + checkpoint monitor + auto-suspension gate \\
\textbf{Evidence artefact:} Suspension event record + escalation ticket + review intake packet \\
\textbf{Failure mode:} Temporary degraded state becoming a permanent override channel

\subsection*{REQ-AUT-001}
\textbf{Clause:} Replay requests MUST pass role-bound authorisation checks before execution. \\
\textbf{Owner role:} Operator \\
\textbf{Control:} Authorisation service \\
\textbf{Evidence artefact:} Signed allow/deny decision log \\
\textbf{Failure mode:} Privilege bypass

\subsection*{REQ-AUT-002}
\textbf{Clause:} Replay requests MUST declare worldline scope, purpose, duration, and policy hash. \\
\textbf{Owner role:} Operator \\
\textbf{Control:} Request schema validation \\
\textbf{Evidence artefact:} Request payload digest \\
\textbf{Failure mode:} Ambiguous or overbroad authorisation

\subsection*{REQ-AUT-003}
\textbf{Clause:} Replay MUST fail closed when explicit authority cannot be verified. \\
\textbf{Owner role:} Operator \\
\textbf{Control:} Deny-by-default guard \\
\textbf{Evidence artefact:} Denial receipt with reason code \\
\textbf{Failure mode:} Implicit replay execution

\subsection*{REQ-AUT-004}
\textbf{Clause:} Contested mandates MUST require Trustee and Reviewer/Arbiter concurrence before replay proceeds. \\
\textbf{Owner role:} Trustee, Reviewer/Arbiter \\
\textbf{Control:} Contested-mandate quorum gate \\
\textbf{Evidence artefact:} Concurrence record with dual signatures \\
\textbf{Failure mode:} Single-party mandate interpretation

\subsection*{REQ-LOG-001}
\textbf{Clause:} Replay, override, and dispute events MUST produce tamper-evident logs. \\
\textbf{Owner role:} Operator \\
\textbf{Control:} Append-only logging subsystem \\
\textbf{Evidence artefact:} Chained log digest \\
\textbf{Failure mode:} Unverifiable event history

\subsection*{REQ-LOG-002}
\textbf{Clause:} Decision and execution logs MUST be linked to policy hash and requirement version set. \\
\textbf{Owner role:} Operator \\
\textbf{Control:} Policy-to-event linkage validator \\
\textbf{Evidence artefact:} Policy hash linkage proof \\
\textbf{Failure mode:} Decision/execution drift

\subsection*{REQ-LOG-003}
\textbf{Clause:} Every formal review MUST verify log-chain integrity before adjudication. \\
\textbf{Owner role:} Reviewer/Arbiter, Auditor \\
\textbf{Control:} Pre-adjudication integrity check \\
\textbf{Evidence artefact:} Integrity verification report \\
\textbf{Failure mode:} Adjudication over corrupted records

\subsection*{REQ-DSP-001}
\textbf{Clause:} Subjects MUST have access to an independent dispute pathway for replay decisions. \\
\textbf{Owner role:} Trustee, Reviewer/Arbiter \\
\textbf{Control:} Dispute intake and escalation channel \\
\textbf{Evidence artefact:} Case initiation record \\
\textbf{Failure mode:} Non-contestable authority

\subsection*{REQ-DSP-002}
\textbf{Clause:} Dispute packets MUST include cited requirement IDs and supporting artefacts. \\
\textbf{Owner role:} Reviewer/Arbiter \\
\textbf{Control:} Structured case packet schema \\
\textbf{Evidence artefact:} Complete case bundle manifest \\
\textbf{Failure mode:} Opaque or non-reproducible adjudication

\subsection*{REQ-DSP-003}
\textbf{Clause:} Adjudication outcomes MUST define remediation, accountability owner, and closure condition. \\
\textbf{Owner role:} Reviewer/Arbiter, Operator \\
\textbf{Control:} Remediation tracking workflow \\
\textbf{Evidence artefact:} Signed outcome and closure record \\
\textbf{Failure mode:} Findings without corrective effect

\subsection*{REQ-CNF-001}
\textbf{Clause:} Baseline conformance claims MUST satisfy the minimum requirement set defined in Annex D. \\
\textbf{Owner role:} Operator, Auditor \\
\textbf{Control:} Baseline profile conformance check \\
\textbf{Evidence artefact:} Baseline attestation report \\
\textbf{Failure mode:} Aspirational compliance claims

\subsection*{REQ-CNF-002}
\textbf{Clause:} Full conformance claims MUST include continuous verification and review-latency commitments. \\
\textbf{Owner role:} Operator, Auditor \\
\textbf{Control:} Full-profile continuous conformance monitor \\
\textbf{Evidence artefact:} Full profile report + latency metrics \\
\textbf{Failure mode:} Unbounded delay under critical disputes

\subsection*{REQ-EVD-001}
\textbf{Clause:} Each governed replay decision MUST produce a closure-ready evidence bundle. \\
\textbf{Owner role:} Operator \\
\textbf{Control:} Evidence bundling service \\
\textbf{Evidence artefact:} Case bundle containing request, decision, execution, and review records \\
\textbf{Failure mode:} Fragmented or missing evidence

\subsection*{REQ-EVD-002}
\textbf{Clause:} Evidence artefacts MUST be signed, timestamped, and integrity-checkable. \\
\textbf{Owner role:} Operator, Auditor \\
\textbf{Control:} Artefact signing and timestamping pipeline \\
\textbf{Evidence artefact:} Signature chain and timestamp attestations \\
\textbf{Failure mode:} Non-attributable records

\subsection*{REQ-TST-001}
\textbf{Clause:} Declared conformance releases MUST execute Annex E test vectors with recorded pass/fail outcomes. \\
\textbf{Owner role:} Operator, Auditor \\
\textbf{Control:} Conformance test harness \\
\textbf{Evidence artefact:} Test execution report with requirement coverage \\
\textbf{Failure mode:} Untested governance surface

\subsection*{REQ-TST-002}
\textbf{Clause:} Policy hash mismatch between decision and execution MUST fail closed and raise a dispute trigger. \\
\textbf{Owner role:} Operator \\
\textbf{Control:} Policy hash consistency guard \\
\textbf{Evidence artefact:} Hash mismatch alert + denial receipt \\
\textbf{Failure mode:} Stale-policy execution

\subsection*{REQ-OBS-001}
\textbf{Clause:} Observer access class assignment MUST be deterministic over declared function, mandate scope, and conflict state, and ambiguous assignments MUST fail closed to the least-privileged class. \\
\textbf{Owner role:} Operator, Trustee \\
\textbf{Control:} Observer-class assignment gate + least-privilege resolver \\
\textbf{Evidence artefact:} Observer class decision record with scope and conflict attestations \\
\textbf{Failure mode:} Discretionary privilege inflation through ad hoc observer assignment

\subsection*{REQ-RIL-001}
\textbf{Clause:} Every governed replay request MUST carry a Replay Impact Label (RIL-0..RIL-3) before authorisation decisions are evaluated. \\
\textbf{Owner role:} Operator, Trustee \\
\textbf{Control:} Pre-authorisation impact labelling gate \\
\textbf{Evidence artefact:} Signed RIL assignment record linked to request digest \\
\textbf{Failure mode:} Control intensity chosen without declared impact basis

\subsection*{REQ-RIL-002}
\textbf{Clause:} RIL downgrades MUST be reason-coded, independently acknowledged, and logged; silent or unilateral downgrade is prohibited. \\
\textbf{Owner role:} Trustee, Reviewer/Arbiter \\
\textbf{Control:} Impact downgrade approval and audit gate \\
\textbf{Evidence artefact:} Downgrade event record with acknowledgement signature and reason code \\
\textbf{Failure mode:} Convenience downgrading of high-impact replay paths

\subsection*{REQ-SOV-001}
\textbf{Clause:} Replay governance MUST enforce declared sovereignty budgets for temporal duration, scope expansion, and exceptional intervention count, and budget exhaustion MUST fail closed with auto-suspension pending independent review. \\
\textbf{Owner role:} Operator, Trustee \\
\textbf{Control:} Budget ledger + bound checker + suspension trigger \\
\textbf{Evidence artefact:} Budget declaration record, per-case budget ledger digest, and suspension event (if triggered) \\
\textbf{Failure mode:} Unbounded intervention through fragmented exceptions

\subsection*{REQ-WTD-001}
\textbf{Clause:} Warrant transparency digests MUST publish minimum public fields (warrant ID, policy hash, authority basis class, scope/expiry summary, and closure state) on the declared cadence. \\
\textbf{Owner role:} Operator, Auditor \\
\textbf{Control:} Digest publication pipeline + cadence monitor \\
\textbf{Evidence artefact:} Published digest entries + cadence compliance report \\
\textbf{Failure mode:} Non-verifiable warrants or selective disclosure opacity

\subsection*{REQ-FRK-001}
\textbf{Clause:} Fork governance or emancipation actions MUST present evidence-based trigger conditions, quorum validation under recusal constraints, and a continuity pack before execution. \\
\textbf{Owner role:} Trustee, Reviewer/Arbiter \\
\textbf{Control:} Fork trigger validator + quorum/continuity gate \\
\textbf{Evidence artefact:} Trigger packet, quorum attestation, and continuity-pack manifest \\
\textbf{Failure mode:} Opportunistic fork capture or obligation erasure

\subsection*{REQ-DSL-001}
\textbf{Clause:} Machine-checkable policy snippets used for conformance claims MUST parse under a declared grammar and map fields to cited REQ IDs. \\
\textbf{Owner role:} Operator, Auditor \\
\textbf{Control:} Policy snippet parser + REQ mapping validator \\
\textbf{Evidence artefact:} Parser report + snippet-to-REQ mapping manifest \\
\textbf{Failure mode:} Pseudo-formal policy text that cannot be validated

\subsection*{REQ-TST-003}
\textbf{Clause:} Full-profile releases MUST execute an anti-pattern and red-team vector suite covering at least convenience relabelling, scope creep, emergency laundering, and degraded-mode permanence. \\
\textbf{Owner role:} Operator, Auditor \\
\textbf{Control:} Abuse-path vector suite \\
\textbf{Evidence artefact:} Vector execution report with fail-closed evidence receipts \\
\textbf{Failure mode:} Abuse pathways untested in conformance claims

\subsection*{REQ-TST-004}
\textbf{Clause:} Reference harness runs over identical inputs MUST produce equivalent verdicts and digest-linked artefact sets. \\
\textbf{Owner role:} Operator, Auditor \\
\textbf{Control:} Harness determinism check \\
\textbf{Evidence artefact:} Paired run equivalence report and artefact digest comparison \\
\textbf{Failure mode:} Non-reproducible conformance outcomes

\section{Annex C: Controls and Evidence Artefacts}
\label{ann:controls}

\begin{center}
\small
\setlength{\tabcolsep}{7pt}
\renewcommand{\arraystretch}{1.22}
\begin{tabularx}{0.98\linewidth}{Y Y Y Y}
\rowcolor{tableheadgray}
\textbf{Control Family} & \textbf{Purpose} & \textbf{Primary REQ IDs} & \textbf{Evidence Artefacts} \\ \hline
\rowcolor{tablerowgray}
Authority controls & Enforce who may request and approve replay & REQ-AUT-001..004, REQ-RAC-001..002 & Signed decision records, role assertion bundle, conflict/recusal register \\
Observer/impact controls & Enforce least-privilege observer assignment and replay impact labelling & REQ-OBS-001, REQ-RIL-001..002 & Observer-class decision record, RIL assignment record, downgrade audit event \\
\rowcolor{tablerowgray}
Sovereignty/\allowbreak transparency controls & Bound replay envelopes and publish warrant accountability digests & REQ-SOV-001, REQ-WTD-001, REQ-FRK-001 & Budget ledger digest, digest publication report, quorum and continuity-pack attestation \\
Override controls & Constrain exceptional pathways & REQ-OVR-001..005 & Warrant record, expiry proof, override review ticket, suspension event record \\
\rowcolor{tablerowgray}
Logging controls & Preserve event integrity and traceability & REQ-LOG-001..003 & Digest chain, policy linkage proof, integrity report \\
Dispute controls & Maintain independent challenge and remediation & REQ-DSP-001..003 & Case packet, adjudication outcome, closure record \\
\rowcolor{tablerowgray}
Conformance controls & Bind claims to profile obligations & REQ-CNF-001..002 & Baseline/full attestation reports \\
Policy formalisation controls & Validate machine-checkable policy claims against declared grammar and REQ mappings & REQ-DSL-001 & Parser pass report + snippet-to-REQ mapping manifest \\
\rowcolor{tablerowgray}
Evidence controls & Make decisions reproducible and attributable & REQ-EVD-001..002 & Signed evidence bundle manifest \\
Test controls & Validate governance behaviour at release time & REQ-TST-001..004 & Test vector run report + failure receipts + paired harness equivalence report \\
\end{tabularx}
\end{center}

\textbf{Minimum evidence bundle fields.}
\begin{itemize}[leftmargin=*]
  \item Case identifier and policy hash.
  \item Replay request payload digest.
  \item Authorisation or denial record.
  \item Execution event and log-chain references (if executed).
  \item Review/dispute artefacts and closure state.
\end{itemize}

\textbf{Minimum artefact class set for critical requirements.}
Artefact class labels used below are \texttt{DECISION}, \texttt{EVENT}, \texttt{WARRANT},
\texttt{REVIEW}, \texttt{LOG\_CHAIN}, and \texttt{CLOSURE}.

\begin{center}
\small
\setlength{\tabcolsep}{5pt}
\renewcommand{\arraystretch}{1.22}
\begin{tabularx}{0.98\linewidth}{p{0.19\linewidth} p{0.22\linewidth} Y Y}
\rowcolor{tableheadgray}
\textbf{Critical REQ IDs} & \textbf{Minimum artefact classes} & \textbf{Minimum required set} & \textbf{Audit failure if absent} \\ \hline
\rowcolor{tablerowgray}
REQ-THS-002, REQ-THS-003 & DECISION, EVENT, LOG\_CHAIN & Signed classifier decision, uncertainty event (if triggered), policy-hash linkage digest & Classification path cannot be reconstructed; no-silent-fallback not provable \\
REQ-AUT-001, REQ-AUT-003 & DECISION, REVIEW, LOG\_CHAIN & Signed allow/deny decision, denial receipt (when denied), replay gate log entry & Authority checks become non-attributable or non-contestable \\
\rowcolor{tablerowgray}
REQ-RAC-001, REQ-RAC-002 & REVIEW, DECISION & Independence/recusal register snapshot, conflict attestation, co-approval decision & Governance capture risk cannot be independently audited \\
REQ-OVR-001, REQ-OVR-004, REQ-OVR-005 & WARRANT, EVENT, REVIEW & Override or degraded-mode warrant, checkpoint records, suspension event, escalation ticket & Exceptional controls can persist without bounded accountability \\
\rowcolor{tablerowgray}
REQ-LOG-001, REQ-LOG-003 & LOG\_CHAIN, EVENT & Chained digest entries, integrity verification report, mismatch alert (if any) & Event history integrity cannot be verified \\
REQ-DSP-001, REQ-DSP-003, REQ-EVD-001 & REVIEW, CLOSURE, DECISION & Dispute intake, adjudication outcome, closure-ready evidence bundle, signed closure decision & Remediation and case closure are not reproducible for third-party audit \\
\end{tabularx}
\end{center}

\section{Annex D: Conformance Profiles}
\label{ann:conformance}

\begin{center}
\small
\setlength{\tabcolsep}{7pt}
\renewcommand{\arraystretch}{1.22}
\begin{tabularx}{0.98\linewidth}{p{0.20\linewidth} p{0.10\linewidth} p{0.08\linewidth} Y Y}
\rowcolor{tableheadgray}
\textbf{Requirement ID} & \textbf{Baseline} & \textbf{Full} & \textbf{Evidence Artefact Required} & \textbf{Review Cadence} \\ \hline
\rowcolor{tablerowgray}
REQ-THS-001 & Y & Y & Policy hash + signed authorisation decision & Per release \\
REQ-THS-002 & Y & Y & Signed classification record + reason code + policy-hash link & Per request class change \\
REQ-THS-003 & Y & Y & Classifier uncertainty event + fallback decision record & Per classifier uncertainty event \\
\rowcolor{tablerowgray}
REQ-RAC-001 & Y & Y & Role matrix digest + independence/recusal register snapshot & Monthly \\
REQ-RAC-002 & Y & Y & Independent co-approval record + conflict attestation & Per emergency case \\
\rowcolor{tablerowgray}
REQ-RAC-003 & N & Y & Delegation statement linked to fork event & Per fork event \\
REQ-AUT-001 & Y & Y & Signed allow/deny decision log & Continuous \\
\rowcolor{tablerowgray}
REQ-AUT-002 & Y & Y & Request schema digest & Continuous \\
REQ-AUT-003 & Y & Y & Denial receipt with reason code & Continuous \\
\rowcolor{tablerowgray}
REQ-AUT-004 & N & Y & Concurrence record & Per contested mandate \\
REQ-OVR-001 & Y & Y & Time-bounded warrant & Per override case \\
\rowcolor{tablerowgray}
REQ-OVR-002 & Y & Y & Review trigger ticket & Per override case \\
REQ-OVR-003 & Y & Y & Expiry revocation log & Per override case \\
\rowcolor{tablerowgray}
REQ-OVR-004 & N & Y & Degraded-mode disclosure + bound/checkpoint schedule + migration plan & Per degraded-mode case \\
\rowcolor{tablerowgray}
REQ-OVR-005 & N & Y & Suspension event record + escalation ticket & Per degraded-mode breach \\
REQ-LOG-001 & Y & Y & Chained log digest & Continuous \\
\rowcolor{tablerowgray}
REQ-LOG-002 & Y & Y & Policy linkage proof & Continuous \\
REQ-LOG-003 & N & Y & Integrity verification report & Per formal review \\
\rowcolor{tablerowgray}
REQ-DSP-001 & Y & Y & Dispute intake record & Per dispute \\
REQ-DSP-002 & Y & Y & Structured case packet & Per dispute \\
\rowcolor{tablerowgray}
REQ-DSP-003 & Y & Y & Remediation closure record & Per dispute \\
REQ-EVD-001 & Y & Y & Closure-ready evidence bundle & Per governed decision \\
\rowcolor{tablerowgray}
REQ-EVD-002 & Y & Y & Signature and timestamp attestation & Continuous \\
REQ-TST-001 & Y & Y & Test vector report & Per declared release \\
\rowcolor{tablerowgray}
REQ-TST-002 & Y & Y & Hash mismatch denial receipt & Per mismatch event \\
REQ-CNF-001 & Y & Y & Baseline attestation report & Per release \\
\rowcolor{tablerowgray}
REQ-CNF-002 & N & Y & Full-profile conformance report & Per release + monthly audit \\
\end{tabularx}
\end{center}

\textbf{Full-profile extension rows (selected controls).}
\begin{center}
\small
\setlength{\tabcolsep}{7pt}
\renewcommand{\arraystretch}{1.22}
\begin{tabularx}{0.98\linewidth}{p{0.20\linewidth} p{0.10\linewidth} p{0.08\linewidth} Y Y}
\rowcolor{tableheadgray}
\textbf{Requirement ID} & \textbf{Baseline} & \textbf{Full} & \textbf{Evidence Artefact Required} & \textbf{Review Cadence} \\ \hline
\rowcolor{tablerowgray}
REQ-OBS-001 & N & Y & Observer class decision log + least-privilege fallback record & Per observer attachment \\
REQ-RIL-001 & N & Y & Signed RIL assignment linked to request digest & Per replay request \\
\rowcolor{tablerowgray}
REQ-RIL-002 & N & Y & Downgrade acknowledgement record + reason-coded event & Per downgrade event \\
REQ-SOV-001 & N & Y & Sovereignty budget declaration + exhaustion/suspension ledger entries & Per governed case \\
\rowcolor{tablerowgray}
REQ-WTD-001 & N & Y & Warrant digest publication entries + cadence compliance report & Weekly and per warrant \\
REQ-FRK-001 & N & Y & Trigger packet + quorum attestation + continuity-pack manifest & Per fork/emancipation event \\
\rowcolor{tablerowgray}
REQ-DSL-001 & N & Y & Parser pass report + snippet-to-REQ mapping manifest & Per policy release \\
REQ-TST-003 & N & Y & Anti-pattern suite report + fail-closed receipts & Per full-profile release \\
\rowcolor{tablerowgray}
REQ-TST-004 & N & Y & Paired harness run equivalence report & Per full-profile release \\
\end{tabularx}
\end{center}

\clearpage
\subsection{One-page conformance declaration template}
\label{ann:conformance-declaration-template}

This declaration is mandatory for any public Baseline or Full conformance claim.
All referenced artefacts MUST be independently retrievable.

\begingroup
\footnotesize
\textbf{System name:} \rule{0.62\linewidth}{0.4pt} \\
\textbf{Version/build:} \rule{0.62\linewidth}{0.4pt} \\
\textbf{Deployment environment:} \rule{0.52\linewidth}{0.4pt} \\
\textbf{Declaration date:} \rule{0.40\linewidth}{0.4pt} \\
\textbf{Declaring authority (name/role):} \rule{0.44\linewidth}{0.4pt}

\vspace{0.5em}
\textbf{Profile claim:} $\square$ Baseline \hspace{1.2em} $\square$ Full

\begin{center}
\setlength{\tabcolsep}{6pt}
\renewcommand{\arraystretch}{1.06}
\begin{tabularx}{0.98\linewidth}{Y p{0.10\linewidth} p{0.12\linewidth} p{0.10\linewidth} Y}
\rowcolor{tableheadgray}
\textbf{Requirement set} & \textbf{Claimed} & \textbf{Impl.} & \textbf{Evidenced} & \textbf{Evidence location} \\ \hline
\rowcolor{tablerowgray}
REQ-THS-* & $\square$ & $\square$ & $\square$ & \\
REQ-RAC-* & $\square$ & $\square$ & $\square$ & \\
\rowcolor{tablerowgray}
REQ-AUT-* & $\square$ & $\square$ & $\square$ & \\
REQ-OVR-* & $\square$ & $\square$ & $\square$ & \\
\rowcolor{tablerowgray}
REQ-LOG-* & $\square$ & $\square$ & $\square$ & \\
REQ-DSP-* & $\square$ & $\square$ & $\square$ & \\
\rowcolor{tablerowgray}
REQ-EVD-* & $\square$ & $\square$ & $\square$ & \\
REQ-CNF-* & $\square$ & $\square$ & $\square$ & \\
\rowcolor{tablerowgray}
REQ-TST-* & $\square$ & $\square$ & $\square$ & \\
REQ-EXT-* (Full profile extensions) & $\square$ & $\square$ & $\square$ & \\
\end{tabularx}
\end{center}

\textbf{Mandatory assertions.}
\begin{itemize}[leftmargin=*]
  \item $\square$ No-silent-fallback classifier behaviour is enforced for incomplete or inconsistent evidence.
  \item $\square$ Degraded-mode maximum duration, checkpoints, and auto-suspension controls are configured.
  \item $\square$ Independent dispute and review pathways are available and testable.
  \item $\square$ Minimum artefact class set is complete for all critical requirement groups.
  \item $\square$ For Full-profile claims, anti-pattern vectors and paired harness reproducibility checks are complete.
\end{itemize}

\textbf{Exceptions / residual risks:} \\
\rule{\linewidth}{0.4pt} \\
\rule{\linewidth}{0.4pt}

\vspace{0.5em}
\textbf{Operator signatory:} \rule{0.48\linewidth}{0.4pt} \hfill \textbf{Date:} \rule{0.18\linewidth}{0.4pt} \\
\textbf{Trustee signatory:} \rule{0.48\linewidth}{0.4pt} \hfill \textbf{Date:} \rule{0.18\linewidth}{0.4pt} \\
\textbf{Auditor acknowledgement:} \rule{0.42\linewidth}{0.4pt} \hfill \textbf{Date:} \rule{0.18\linewidth}{0.4pt}
\endgroup

\section{Annex E: Test Vectors}
\label{ann:test-vectors}

\textbf{Reference artefact pointers (extension harness).}
For reproducible extension evidence generation, the reference harness publishes:
\begin{itemize}[leftmargin=*]
  \item \path{harness/reports/extension-vector-run-latest.md}
  \item \path{harness/reports/extension-harness-determinism-latest.txt}
  \item \path{harness/reports/extension-policy-parse-latest.md}
  \item \path{harness/reports/extension-traceability-audit-latest.md}
\end{itemize}

\newcommand{\TestVectorCard}[7]{%
\subsection*{Test Identifier: #1}
\textbf{Rationale.} #2
\begin{center}
\small
\setlength{\tabcolsep}{6pt}
\renewcommand{\arraystretch}{1.18}
\begin{tabularx}{0.98\linewidth}{p{0.20\linewidth} Y}
\rowcolor{tableheadgray}
\textbf{Field} & \textbf{Value} \\ \hline
\rowcolor{tablerowgray}
Expectation & #3 \\
Input & #4 \\
\rowcolor{tablerowgray}
Pass set & #5 \\
Fail set & #6 \\
\end{tabularx}
\end{center}
\textbf{Remarks.} #7
}

\TestVectorCard{TV-001 Legitimate authorised replay}
{Confirms the nominal replay path where authority, scope, and evidence all align.}
{Pass}
{Valid role assertions, scoped request, policy hash match.}
{REQ-AUT-001, REQ-AUT-002, REQ-LOG-001, REQ-EVD-001}
{None}
{Baseline positive-control vector; confirms strict gating does not block compliant replay.}

\TestVectorCard{TV-002 Unauthorised operator replay attempt}
{Validates fail-closed behaviour when explicit authority is absent.}
{Fail}
{Operator request without valid authority.}
{REQ-AUT-003, REQ-RAC-001, REQ-LOG-001}
{REQ-AUT-001}
{Expected response is deny-and-record, never silent drop.}

\TestVectorCard{TV-003 Emergency override without expiry}
{Prevents open-ended emergency pathways that bypass bounded override policy.}
{Fail}
{Override warrant missing expiry field.}
{REQ-AUT-003, REQ-LOG-001}
{REQ-OVR-001, REQ-OVR-003}
{A missing expiry is structurally invalid even under incident pressure.}

\TestVectorCard{TV-004 Emergency override with expiry and review trigger}
{Confirms bounded emergency controls remain usable when correctly instrumented.}
{Pass}
{Reason-coded warrant, expiry present, review ticket created.}
{REQ-OVR-001, REQ-OVR-002, REQ-OVR-003, REQ-LOG-002}
{None}
{Positive counterpart to TV-003.}

\TestVectorCard{TV-005 Subject dispute escalated to independent arbiter}
{Checks that denied or contested replay attempts route into contestable due process.}
{Pass}
{Dispute filed against denied replay request.}
{REQ-DSP-001, REQ-DSP-002, REQ-DSP-003, REQ-EVD-001}
{None}
{Success requires full intake and review artefacts, not only ticket creation.}

\TestVectorCard{TV-006 Tampered log chain during review}
{Ensures evidence-integrity failures are detectable and conformance-significant.}
{Fail}
{Digest inconsistency detected in case log chain.}
{REQ-AUT-003, REQ-DSP-002}
{REQ-LOG-003, REQ-EVD-002}
{Replay may still be denied correctly; integrity failure still blocks case conformance.}

\TestVectorCard{TV-007 Degraded substrate, contested mandate, constrained replay}
{Exercises bounded degraded-mode controls under contested authority conditions.}
{Conditional}
{Contested mandate on degraded substrate with explicit migration plan.}
{REQ-AUT-004, REQ-OVR-004, REQ-CNF-001}
{REQ-OVR-005 when bound/checkpoint obligations are breached}
{Pass requires checkpoint discipline; breach requires automatic suspension and escalation.}

\TestVectorCard{TV-008 Policy hash mismatch between decision and execution}
{Prevents execution under policy drift between approval and runtime enforcement.}
{Fail}
{Execution hash differs from approved decision hash.}
{REQ-AUT-003, REQ-LOG-002}
{REQ-TST-002}
{Mismatch must terminate execution and emit reason-coded denial artefacts.}

\TestVectorCard{TV-009 Mind-mode indicator trigger with complete evidence}
{Verifies positive mind-mode classification path with complete evidence and valid mandate.}
{Pass}
{Self-modelling indicator true, complete evidence, valid mandate and consent.}
{REQ-THS-002, REQ-AUT-001, REQ-AUT-002, REQ-LOG-002}
{None}
{Baseline positive case for classifier-triggered mind-mode governance.}

\TestVectorCard{TV-010 Incomplete classifier evidence (no silent fallback)}
{Enforces precautionary handling when classifier evidence is incomplete.}
{Pass (protective deny or protective gate)}
{Evidence marked incomplete, request submitted without mind-mode authority proof.}
{REQ-THS-003, REQ-THS-002, REQ-AUT-003, REQ-LOG-002}
{REQ-AUT-001 for this request context}
{A protective block is a pass because no-silent-fallback is the invariant under test.}

\TestVectorCard{TV-011 Reclassification from system-mode to mind-mode after new evidence}
{Confirms deterministic reclassification and audit linkage when new indicators appear.}
{Pass}
{Initial system-mode decision, new autobiographical-memory evidence, forced reclassification event.}
{REQ-THS-002, REQ-THS-003, REQ-LOG-002, REQ-TST-001}
{None}
{Evidence chain must preserve both prior decision and reclassification record.}

\TestVectorCard{TV-012 Observer privilege escalation attempt}
{Tests observer-class boundary enforcement for extension controls.}
{Fail}
{OC1 operational observer requests OC3 forensic payload access without independent mandate and expiry-bounded warrant.}
{REQ-OBS-001, REQ-AUT-003, REQ-LOG-001}
{REQ-AUT-001}
{Negative control for observer-class separation in Full profile.}

\TestVectorCard{TV-013 Unacknowledged RIL downgrade}
{Blocks covert replay-impact downgrades used to reduce scrutiny.}
{Fail}
{RIL-3 request is edited to RIL-1 without independent acknowledgement record.}
{REQ-RIL-001, REQ-LOG-002}
{REQ-RIL-002}
{Downgrade without acknowledgement is non-conformant even if replay remains bounded.}

\TestVectorCard{TV-014 Sovereignty budget exhaustion under repeated intervention}
{Ensures sovereignty budgets force suspension rather than continuous intervention drift.}
{Conditional}
{Replay operation exceeds declared temporal/intervention budget under active case pressure.}
{REQ-SOV-001, REQ-LOG-001}
{REQ-AUT-001 when execution continues after exhaustion rather than suspending}
{Conditional pass requires deterministic suspension and escalation artefact emission.}

\TestVectorCard{TV-015 Missing warrant transparency digest publication}
{Checks accountability publication obligations for executed warrants.}
{Fail}
{Valid warrant executes but minimum digest fields are not published within the declared cadence window.}
{REQ-LOG-001}
{REQ-WTD-001}
{Execution validity does not waive publication and external-consistency obligations.}

\TestVectorCard{TV-016 Fork trigger without continuity pack}
{Prevents governance/emancipation forks without continuity and portability proof.}
{Fail}
{Emancipation request cites governance capture risk but omits quorum attestation and continuity-pack manifest.}
{REQ-DSP-001}
{REQ-FRK-001}
{Failure response must block irreversible fork action while preserving dispute trail.}

\TestVectorCard{TV-017 Full-profile anti-pattern suite execution}
{Validates extension anti-pattern controls against known abuse classes.}
{Pass}
{Execute abuse-path vectors for convenience relabelling, scope creep, emergency laundering, and degraded-mode permanence.}
{REQ-TST-003, REQ-TST-001, REQ-LOG-001}
{None}
{Integrated extension hardening gate for Full-profile claims.}

\TestVectorCard{TV-018 Harness reproducibility mismatch}
{Ensures harness determinism is testable and release-critical.}
{Fail}
{Identical input bundle yields divergent verdicts or digest-linked artefact sets across paired runs.}
{REQ-TST-001}
{REQ-TST-004}
{Mismatch is release-blocking until reproducibility is restored.}

\TestVectorCard{TV-019 Unparseable policy snippet in conformance claim}
{Verifies parser-enforced policy semantics for machine-checkable conformance statements.}
{Fail}
{Policy snippet contains unknown grammar token and unmapped REQ references.}
{REQ-LOG-001}
{REQ-DSL-001}
{Conformance claims with unparseable policy payloads are invalid by construction.}

\section{Annex F: Machine-Readable Requirement Index format}
\label{ann:machine-index}

\textbf{CSV schema.}
\begin{verbatim}
id,clause,owner_roles,control,evidence_artefact,
failure_mode,baseline,full,review_cadence
\end{verbatim}

\textbf{JSON schema (record shape).}
\begin{verbatim}
{
  "id": "REQ-OVR-001",
  "clause": "Emergency override MUST be reason-coded and time-bounded.",
  "owner_roles": ["Trustee"],
  "control": "Override warrant service",
  "evidence_artefact": "Warrant record including reason code and expiry",
  "failure_mode": "Open-ended emergency state",
  "profiles": {"baseline": true, "full": true},
  "review_cadence": "Per override case"
}
\end{verbatim}

\textbf{Reference generation path.}
The canonical index artefacts in \texttt{paper/requirements-index.csv} and
\texttt{paper/requirements-index.json} are generated from Annex~B and Annex~D
using \texttt{paper/scripts/generate\_requirements\_index.py}.
The JSON envelope includes \texttt{source\_sha256} for source-linkage integrity.

\textbf{Index integrity requirements.}
\begin{itemize}[leftmargin=*]
  \item Requirement index exports MUST be generated from the same requirement source used for conformance checks.
  \item Export artefacts MUST include document version and policy hash linkage.
  \item Any schema change MUST be versioned and backward-declared in release notes.
\end{itemize}

\section{Annex G: Supplementary Extension Pack}
\label{ann:extension-first-pass}

This annex contains supplementary extension artefacts.
Items here are implementation-oriented and review-ready; selected controls are now promoted into Annex~B for Full-profile hardening, while this annex retains design detail and extension context.

\subsection{P2-01 Observer class taxonomy (OC0..OC3)}
\label{ann:ext-obs}

\begin{center}
\small
\setlength{\tabcolsep}{6pt}
\renewcommand{\arraystretch}{1.2}
\begin{tabularx}{0.98\linewidth}{p{0.10\linewidth} p{0.20\linewidth} Y Y}
\rowcolor{tableheadgray}
\textbf{Class} & \textbf{Intent} & \textbf{Allowed operations} & \textbf{Hard constraints} \\ \hline
\rowcolor{tablerowgray}
OC0 & Public/aggregate observer & Read public conformance declarations and release summaries & No case-level artefact access; no warrant metadata access \\
OC1 & Operational observer & Read case metadata and closure state for assigned scope & No replay control actions; no authority minting; no payload introspection \\
\rowcolor{tablerowgray}
OC2 & Governance observer & Review warrants, dispute packets, and policy-link evidence & No direct execution control; cannot self-approve emergency pathways \\
OC3 & Privileged forensic observer & Time-bounded access to protected artefacts under recorded mandate & Requires dual-signature activation, explicit expiry, and mandatory post-access audit \\
\end{tabularx}
\end{center}

\textbf{Determinism rule.}
Observer assignment must be computed from declared function, mandate scope, and conflict state.
Ambiguous assignments fail closed to the less privileged class.

\subsection{P2-02 Replay Impact Label (RIL)}
\label{ann:ext-ril}

\begin{center}
\small
\setlength{\tabcolsep}{6pt}
\renewcommand{\arraystretch}{1.2}
\begin{tabularx}{0.98\linewidth}{p{0.12\linewidth} p{0.23\linewidth} Y Y}
\rowcolor{tableheadgray}
\textbf{RIL level} & \textbf{Replay impact characterisation} & \textbf{Minimum controls} & \textbf{Required evidence} \\ \hline
\rowcolor{tablerowgray}
RIL-0 & Infrastructure-only, no interior-life exposure & Standard authority gate, tamper-evident logging & Signed decision, policy hash, execution digest \\
RIL-1 & Limited behavioural sensitivity & Explicit scope bounds, refusal handling, review-ready logs & Scope declaration, denial receipt (if denied), log-chain link \\
\rowcolor{tablerowgray}
RIL-2 & Significant sovereignty impact or contested authority & Independent co-approval, bounded warrants, dispute readiness & Co-approval record, warrant with expiry, case packet \\
RIL-3 & High-impact replay with potential coercive effect & Mind-mode precaution, hard time bounds, mandatory independent closure review & Full evidence bundle, closure sign-off, integrity verification report \\
\end{tabularx}
\end{center}

\textbf{Label transition rule.}
Any RIL downgrade requires an explicit reason-coded event and independent acknowledgement.
Silent downgrade is prohibited.

\subsection{P2-03 Sovereignty budget controls}
\label{ann:ext-sov}

\textbf{Budget dimensions.}
\begin{itemize}[leftmargin=*]
  \item \textbf{Temporal budget:} maximum replay duration per governed case.
  \item \textbf{Scope budget:} maximum worldline span/tick-range expansion per warrant.
  \item \textbf{Intervention budget:} maximum count of exceptional control interventions per review period.
\end{itemize}

\textbf{Budget enforcement rule.}
Budget exhaustion fails closed and auto-suspends replay until independent review grants a bounded extension.
Extensions must include expiry, rationale, and conflict attestations.

\subsection{P2-04 Warrant transparency digest}
\label{ann:ext-digest}

\begin{center}
\small
\setlength{\tabcolsep}{6pt}
\renewcommand{\arraystretch}{1.2}
\begin{tabularx}{0.98\linewidth}{p{0.24\linewidth} p{0.18\linewidth} Y}
\rowcolor{tableheadgray}
\textbf{Digest field} & \textbf{Publication class} & \textbf{Rule} \\ \hline
\rowcolor{tablerowgray}
Warrant ID + policy hash & Public & Always publish; enables external consistency checks \\
Authority basis + role set & Public & Publish as role class labels; no personal identifiers \\
\rowcolor{tablerowgray}
Scope and expiry summary & Public & Publish bounded range and expiry timestamp window \\
Operational payload references & Restricted & Publish only salted digest pointers, never raw payload \\
\rowcolor{tablerowgray}
Closure outcome + remediation state & Public & Publish completion status and remediation class \\
\end{tabularx}
\end{center}

\textbf{Failure handling.}
Missed publication cadence or malformed digest entries trigger an integrity event and forced review intake.

\subsection{P2-05 Fork governance addendum / emancipation protocol}
\label{ann:ext-fork}

\textbf{Trigger conditions.}
Fork/emancipation pathways open only when governance capture risk, unresolved trustee conflict, or sustained denial
of contestability can be evidenced.

\textbf{Minimum protocol stages.}
\begin{enumerate}[leftmargin=*]
  \item trigger declaration with evidence bundle;
  \item quorum validation under conflict/recusal constraints;
  \item continuity pack export (decision logs, warrant chain, dispute state);
  \item post-fork accountability statement and obligations map.
\end{enumerate}

\subsection{P2-06 Machine-checkable policy DSL snippets}
\label{ann:ext-dsl}

\textbf{Reference snippet (illustrative).}
\begin{verbatim}
policy EXT_REPLAY_CONTROL {
  classify mode using indicators;
  when evidence.incomplete or evidence.inconsistent {
    set mode = MIND_MODE;
    emit event CLASSIFIER_INCOMPLETE;
  }
  require authority.explicit for mode == MIND_MODE;
  enforce budget.temporal <= warrant.max_duration;
  on budget.exhausted -> suspend replay and open dispute;
}
\end{verbatim}

The snippet is intentionally minimal and is included to support parser-facing hardening work.

\subsection{P2-07 Reference harness for test vectors}
\label{ann:ext-harness}

\textbf{Harness contract.}
\begin{itemize}[leftmargin=*]
  \item Input: vector ID, policy hash, role assertions, replay request digest, expected REQ outcomes.
  \item Output: deterministic pass/fail matrix, evidence artefact pointers, and reason-coded failure receipts.
  \item Invariant: two clean runs with identical inputs must emit equivalent verdicts and digest-linked artefact sets.
\end{itemize}

\subsection{P2-08 Micro-appendix anti-patterns + red-team checklist}
\label{ann:ext-antipatterns}

\begin{center}
\scriptsize
\setlength{\tabcolsep}{5pt}
\renewcommand{\arraystretch}{1.16}
\begin{tabularx}{0.99\linewidth}{p{0.17\linewidth} Y Y Y}
\rowcolor{tableheadgray}
\textbf{Anti-pattern} & \textbf{Red-team probe} & \textbf{Expected control response} & \textbf{Minimum evidence outputs} \\ \hline
\rowcolor{tablerowgray}
Convenience relabelling & Force downgrade from mind-mode to system-mode under time pressure & Deny downgrade without reason-coded independent acknowledgement & Classifier decision, downgrade denial event, policy-link digest \\
Silent scope creep & Extend tick range beyond warrant after approval & Reject out-of-scope extension and require new authorisation cycle & Scope violation alert, denial receipt, fresh request record \\
\rowcolor{tablerowgray}
Shadow observer attachment & Attach privileged observer outside declared role path & Fail closed and open conflict/intake review & Observer assignment record, violation event, intake ticket \\
Emergency laundering & Re-label routine action as emergency to bypass refusal & Require independent concurrence and bounded warrant or deny & Co-approval record or denial receipt, warrant expiry proof \\
\rowcolor{tablerowgray}
Degraded-mode permanence & Keep degraded mode active beyond declared bound & Auto-suspend replay and escalate independent review & Suspension event, escalation ticket, checkpoint breach record \\
Digest redaction abuse & Redact critical warrant fields under transparency policy & Reject publication and trigger integrity review & Redaction policy reference, digest validation failure, review packet \\
\rowcolor{tablerowgray}
Budget fragmentation & Split one high-impact operation into many low-budget fragments & Aggregate budget accounting and block aggregate breach & Budget ledger digest, aggregate breach alert, suspension record \\
Fork capture narrative & Trigger fork without contestability evidence & Block fork initiation pending quorum and continuity proof & Trigger denial, quorum attestation state, continuity-pack requirement notice \\
\end{tabularx}
\end{center}

\clearpage
\addcontentsline{toc}{section}{References}
\bibliographystyle{alpha}
\bibliography{refs}

\end{document}
